\documentclass[11pt]{article}

    \usepackage[breakable]{tcolorbox}
    \usepackage{parskip} % Stop auto-indenting (to mimic markdown behaviour)
    

    % Basic figure setup, for now with no caption control since it's done
    % automatically by Pandoc (which extracts ![](path) syntax from Markdown).
    \usepackage{graphicx}
    % Keep aspect ratio if custom image width or height is specified
    \setkeys{Gin}{keepaspectratio}
    % Maintain compatibility with old templates. Remove in nbconvert 6.0
    \let\Oldincludegraphics\includegraphics
    % Ensure that by default, figures have no caption (until we provide a
    % proper Figure object with a Caption API and a way to capture that
    % in the conversion process - todo).
    \usepackage{caption}
    \DeclareCaptionFormat{nocaption}{}
    \captionsetup{format=nocaption,aboveskip=0pt,belowskip=0pt}

    \usepackage{float}
    \floatplacement{figure}{H} % forces figures to be placed at the correct location
    \usepackage{xcolor} % Allow colors to be defined
    \usepackage{enumerate} % Needed for markdown enumerations to work
    \usepackage{geometry} % Used to adjust the document margins
    \usepackage{amsmath} % Equations
    \usepackage{amssymb} % Equations
    \usepackage{textcomp} % defines textquotesingle
    % Hack from http://tex.stackexchange.com/a/47451/13684:
    \AtBeginDocument{%
        \def\PYZsq{\textquotesingle}% Upright quotes in Pygmentized code
    }
    \usepackage{upquote} % Upright quotes for verbatim code
    \usepackage{eurosym} % defines \euro

    \usepackage{iftex}
    \ifPDFTeX
        \usepackage[T1]{fontenc}
        \IfFileExists{alphabeta.sty}{
              \usepackage{alphabeta}
          }{
              \usepackage[mathletters]{ucs}
              \usepackage[utf8x]{inputenc}
          }
    \else
        \usepackage{fontspec}
        \usepackage{unicode-math}
    \fi

    \usepackage{fancyvrb} % verbatim replacement that allows latex
    \usepackage{grffile} % extends the file name processing of package graphics
                         % to support a larger range
    \makeatletter % fix for old versions of grffile with XeLaTeX
    \@ifpackagelater{grffile}{2019/11/01}
    {
      % Do nothing on new versions
    }
    {
      \def\Gread@@xetex#1{%
        \IfFileExists{"\Gin@base".bb}%
        {\Gread@eps{\Gin@base.bb}}%
        {\Gread@@xetex@aux#1}%
      }
    }
    \makeatother
    \usepackage[Export]{adjustbox} % Used to constrain images to a maximum size
    \adjustboxset{max size={0.9\linewidth}{0.9\paperheight}}

    % The hyperref package gives us a pdf with properly built
    % internal navigation ('pdf bookmarks' for the table of contents,
    % internal cross-reference links, web links for URLs, etc.)
    \usepackage{hyperref}
    % The default LaTeX title has an obnoxious amount of whitespace. By default,
    % titling removes some of it. It also provides customization options.
    \usepackage{titling}
    \usepackage{longtable} % longtable support required by pandoc >1.10
    \usepackage{booktabs}  % table support for pandoc > 1.12.2
    \usepackage{array}     % table support for pandoc >= 2.11.3
    \usepackage{calc}      % table minipage width calculation for pandoc >= 2.11.1
    \usepackage[inline]{enumitem} % IRkernel/repr support (it uses the enumerate* environment)
    \usepackage[normalem]{ulem} % ulem is needed to support strikethroughs (\sout)
                                % normalem makes italics be italics, not underlines
    \usepackage{soul}      % strikethrough (\st) support for pandoc >= 3.0.0
    \usepackage{mathrsfs}
    

    
    % Colors for the hyperref package
    \definecolor{urlcolor}{rgb}{0,.145,.698}
    \definecolor{linkcolor}{rgb}{.71,0.21,0.01}
    \definecolor{citecolor}{rgb}{.12,.54,.11}

    % ANSI colors
    \definecolor{ansi-black}{HTML}{3E424D}
    \definecolor{ansi-black-intense}{HTML}{282C36}
    \definecolor{ansi-red}{HTML}{E75C58}
    \definecolor{ansi-red-intense}{HTML}{B22B31}
    \definecolor{ansi-green}{HTML}{00A250}
    \definecolor{ansi-green-intense}{HTML}{007427}
    \definecolor{ansi-yellow}{HTML}{DDB62B}
    \definecolor{ansi-yellow-intense}{HTML}{B27D12}
    \definecolor{ansi-blue}{HTML}{208FFB}
    \definecolor{ansi-blue-intense}{HTML}{0065CA}
    \definecolor{ansi-magenta}{HTML}{D160C4}
    \definecolor{ansi-magenta-intense}{HTML}{A03196}
    \definecolor{ansi-cyan}{HTML}{60C6C8}
    \definecolor{ansi-cyan-intense}{HTML}{258F8F}
    \definecolor{ansi-white}{HTML}{C5C1B4}
    \definecolor{ansi-white-intense}{HTML}{A1A6B2}
    \definecolor{ansi-default-inverse-fg}{HTML}{FFFFFF}
    \definecolor{ansi-default-inverse-bg}{HTML}{000000}

    % common color for the border for error outputs.
    \definecolor{outerrorbackground}{HTML}{FFDFDF}

    % commands and environments needed by pandoc snippets
    % extracted from the output of `pandoc -s`
    \providecommand{\tightlist}{%
      \setlength{\itemsep}{0pt}\setlength{\parskip}{0pt}}
    \DefineVerbatimEnvironment{Highlighting}{Verbatim}{commandchars=\\\{\}}
    % Add ',fontsize=\small' for more characters per line
    \newenvironment{Shaded}{}{}
    \newcommand{\KeywordTok}[1]{\textcolor[rgb]{0.00,0.44,0.13}{\textbf{{#1}}}}
    \newcommand{\DataTypeTok}[1]{\textcolor[rgb]{0.56,0.13,0.00}{{#1}}}
    \newcommand{\DecValTok}[1]{\textcolor[rgb]{0.25,0.63,0.44}{{#1}}}
    \newcommand{\BaseNTok}[1]{\textcolor[rgb]{0.25,0.63,0.44}{{#1}}}
    \newcommand{\FloatTok}[1]{\textcolor[rgb]{0.25,0.63,0.44}{{#1}}}
    \newcommand{\CharTok}[1]{\textcolor[rgb]{0.25,0.44,0.63}{{#1}}}
    \newcommand{\StringTok}[1]{\textcolor[rgb]{0.25,0.44,0.63}{{#1}}}
    \newcommand{\CommentTok}[1]{\textcolor[rgb]{0.38,0.63,0.69}{\textit{{#1}}}}
    \newcommand{\OtherTok}[1]{\textcolor[rgb]{0.00,0.44,0.13}{{#1}}}
    \newcommand{\AlertTok}[1]{\textcolor[rgb]{1.00,0.00,0.00}{\textbf{{#1}}}}
    \newcommand{\FunctionTok}[1]{\textcolor[rgb]{0.02,0.16,0.49}{{#1}}}
    \newcommand{\RegionMarkerTok}[1]{{#1}}
    \newcommand{\ErrorTok}[1]{\textcolor[rgb]{1.00,0.00,0.00}{\textbf{{#1}}}}
    \newcommand{\NormalTok}[1]{{#1}}

    % Additional commands for more recent versions of Pandoc
    \newcommand{\ConstantTok}[1]{\textcolor[rgb]{0.53,0.00,0.00}{{#1}}}
    \newcommand{\SpecialCharTok}[1]{\textcolor[rgb]{0.25,0.44,0.63}{{#1}}}
    \newcommand{\VerbatimStringTok}[1]{\textcolor[rgb]{0.25,0.44,0.63}{{#1}}}
    \newcommand{\SpecialStringTok}[1]{\textcolor[rgb]{0.73,0.40,0.53}{{#1}}}
    \newcommand{\ImportTok}[1]{{#1}}
    \newcommand{\DocumentationTok}[1]{\textcolor[rgb]{0.73,0.13,0.13}{\textit{{#1}}}}
    \newcommand{\AnnotationTok}[1]{\textcolor[rgb]{0.38,0.63,0.69}{\textbf{\textit{{#1}}}}}
    \newcommand{\CommentVarTok}[1]{\textcolor[rgb]{0.38,0.63,0.69}{\textbf{\textit{{#1}}}}}
    \newcommand{\VariableTok}[1]{\textcolor[rgb]{0.10,0.09,0.49}{{#1}}}
    \newcommand{\ControlFlowTok}[1]{\textcolor[rgb]{0.00,0.44,0.13}{\textbf{{#1}}}}
    \newcommand{\OperatorTok}[1]{\textcolor[rgb]{0.40,0.40,0.40}{{#1}}}
    \newcommand{\BuiltInTok}[1]{{#1}}
    \newcommand{\ExtensionTok}[1]{{#1}}
    \newcommand{\PreprocessorTok}[1]{\textcolor[rgb]{0.74,0.48,0.00}{{#1}}}
    \newcommand{\AttributeTok}[1]{\textcolor[rgb]{0.49,0.56,0.16}{{#1}}}
    \newcommand{\InformationTok}[1]{\textcolor[rgb]{0.38,0.63,0.69}{\textbf{\textit{{#1}}}}}
    \newcommand{\WarningTok}[1]{\textcolor[rgb]{0.38,0.63,0.69}{\textbf{\textit{{#1}}}}}
    \makeatletter
    \newsavebox\pandoc@box
    \newcommand*\pandocbounded[1]{%
      \sbox\pandoc@box{#1}%
      % scaling factors for width and height
      \Gscale@div\@tempa\textheight{\dimexpr\ht\pandoc@box+\dp\pandoc@box\relax}%
      \Gscale@div\@tempb\linewidth{\wd\pandoc@box}%
      % select the smaller of both
      \ifdim\@tempb\p@<\@tempa\p@
        \let\@tempa\@tempb
      \fi
      % scaling accordingly (\@tempa < 1)
      \ifdim\@tempa\p@<\p@
        \scalebox{\@tempa}{\usebox\pandoc@box}%
      % scaling not needed, use as it is
      \else
        \usebox{\pandoc@box}%
      \fi
    }
    \makeatother

    % Define a nice break command that doesn't care if a line doesn't already
    % exist.
    \def\br{\hspace*{\fill} \\* }
    % Math Jax compatibility definitions
    \def\gt{>}
    \def\lt{<}
    \let\Oldtex\TeX
    \let\Oldlatex\LaTeX
    \renewcommand{\TeX}{\textrm{\Oldtex}}
    \renewcommand{\LaTeX}{\textrm{\Oldlatex}}
    % Document parameters
    % Document title
    \title{讲义\_无代码}
    
    
    
    
    
    
    
% Pygments definitions
\makeatletter
\def\PY@reset{\let\PY@it=\relax \let\PY@bf=\relax%
    \let\PY@ul=\relax \let\PY@tc=\relax%
    \let\PY@bc=\relax \let\PY@ff=\relax}
\def\PY@tok#1{\csname PY@tok@#1\endcsname}
\def\PY@toks#1+{\ifx\relax#1\empty\else%
    \PY@tok{#1}\expandafter\PY@toks\fi}
\def\PY@do#1{\PY@bc{\PY@tc{\PY@ul{%
    \PY@it{\PY@bf{\PY@ff{#1}}}}}}}
\def\PY#1#2{\PY@reset\PY@toks#1+\relax+\PY@do{#2}}

\@namedef{PY@tok@w}{\def\PY@tc##1{\textcolor[rgb]{0.73,0.73,0.73}{##1}}}
\@namedef{PY@tok@c}{\let\PY@it=\textit\def\PY@tc##1{\textcolor[rgb]{0.24,0.48,0.48}{##1}}}
\@namedef{PY@tok@cp}{\def\PY@tc##1{\textcolor[rgb]{0.61,0.40,0.00}{##1}}}
\@namedef{PY@tok@k}{\let\PY@bf=\textbf\def\PY@tc##1{\textcolor[rgb]{0.00,0.50,0.00}{##1}}}
\@namedef{PY@tok@kp}{\def\PY@tc##1{\textcolor[rgb]{0.00,0.50,0.00}{##1}}}
\@namedef{PY@tok@kt}{\def\PY@tc##1{\textcolor[rgb]{0.69,0.00,0.25}{##1}}}
\@namedef{PY@tok@o}{\def\PY@tc##1{\textcolor[rgb]{0.40,0.40,0.40}{##1}}}
\@namedef{PY@tok@ow}{\let\PY@bf=\textbf\def\PY@tc##1{\textcolor[rgb]{0.67,0.13,1.00}{##1}}}
\@namedef{PY@tok@nb}{\def\PY@tc##1{\textcolor[rgb]{0.00,0.50,0.00}{##1}}}
\@namedef{PY@tok@nf}{\def\PY@tc##1{\textcolor[rgb]{0.00,0.00,1.00}{##1}}}
\@namedef{PY@tok@nc}{\let\PY@bf=\textbf\def\PY@tc##1{\textcolor[rgb]{0.00,0.00,1.00}{##1}}}
\@namedef{PY@tok@nn}{\let\PY@bf=\textbf\def\PY@tc##1{\textcolor[rgb]{0.00,0.00,1.00}{##1}}}
\@namedef{PY@tok@ne}{\let\PY@bf=\textbf\def\PY@tc##1{\textcolor[rgb]{0.80,0.25,0.22}{##1}}}
\@namedef{PY@tok@nv}{\def\PY@tc##1{\textcolor[rgb]{0.10,0.09,0.49}{##1}}}
\@namedef{PY@tok@no}{\def\PY@tc##1{\textcolor[rgb]{0.53,0.00,0.00}{##1}}}
\@namedef{PY@tok@nl}{\def\PY@tc##1{\textcolor[rgb]{0.46,0.46,0.00}{##1}}}
\@namedef{PY@tok@ni}{\let\PY@bf=\textbf\def\PY@tc##1{\textcolor[rgb]{0.44,0.44,0.44}{##1}}}
\@namedef{PY@tok@na}{\def\PY@tc##1{\textcolor[rgb]{0.41,0.47,0.13}{##1}}}
\@namedef{PY@tok@nt}{\let\PY@bf=\textbf\def\PY@tc##1{\textcolor[rgb]{0.00,0.50,0.00}{##1}}}
\@namedef{PY@tok@nd}{\def\PY@tc##1{\textcolor[rgb]{0.67,0.13,1.00}{##1}}}
\@namedef{PY@tok@s}{\def\PY@tc##1{\textcolor[rgb]{0.73,0.13,0.13}{##1}}}
\@namedef{PY@tok@sd}{\let\PY@it=\textit\def\PY@tc##1{\textcolor[rgb]{0.73,0.13,0.13}{##1}}}
\@namedef{PY@tok@si}{\let\PY@bf=\textbf\def\PY@tc##1{\textcolor[rgb]{0.64,0.35,0.47}{##1}}}
\@namedef{PY@tok@se}{\let\PY@bf=\textbf\def\PY@tc##1{\textcolor[rgb]{0.67,0.36,0.12}{##1}}}
\@namedef{PY@tok@sr}{\def\PY@tc##1{\textcolor[rgb]{0.64,0.35,0.47}{##1}}}
\@namedef{PY@tok@ss}{\def\PY@tc##1{\textcolor[rgb]{0.10,0.09,0.49}{##1}}}
\@namedef{PY@tok@sx}{\def\PY@tc##1{\textcolor[rgb]{0.00,0.50,0.00}{##1}}}
\@namedef{PY@tok@m}{\def\PY@tc##1{\textcolor[rgb]{0.40,0.40,0.40}{##1}}}
\@namedef{PY@tok@gh}{\let\PY@bf=\textbf\def\PY@tc##1{\textcolor[rgb]{0.00,0.00,0.50}{##1}}}
\@namedef{PY@tok@gu}{\let\PY@bf=\textbf\def\PY@tc##1{\textcolor[rgb]{0.50,0.00,0.50}{##1}}}
\@namedef{PY@tok@gd}{\def\PY@tc##1{\textcolor[rgb]{0.63,0.00,0.00}{##1}}}
\@namedef{PY@tok@gi}{\def\PY@tc##1{\textcolor[rgb]{0.00,0.52,0.00}{##1}}}
\@namedef{PY@tok@gr}{\def\PY@tc##1{\textcolor[rgb]{0.89,0.00,0.00}{##1}}}
\@namedef{PY@tok@ge}{\let\PY@it=\textit}
\@namedef{PY@tok@gs}{\let\PY@bf=\textbf}
\@namedef{PY@tok@ges}{\let\PY@bf=\textbf\let\PY@it=\textit}
\@namedef{PY@tok@gp}{\let\PY@bf=\textbf\def\PY@tc##1{\textcolor[rgb]{0.00,0.00,0.50}{##1}}}
\@namedef{PY@tok@go}{\def\PY@tc##1{\textcolor[rgb]{0.44,0.44,0.44}{##1}}}
\@namedef{PY@tok@gt}{\def\PY@tc##1{\textcolor[rgb]{0.00,0.27,0.87}{##1}}}
\@namedef{PY@tok@err}{\def\PY@bc##1{{\setlength{\fboxsep}{\string -\fboxrule}\fcolorbox[rgb]{1.00,0.00,0.00}{1,1,1}{\strut ##1}}}}
\@namedef{PY@tok@kc}{\let\PY@bf=\textbf\def\PY@tc##1{\textcolor[rgb]{0.00,0.50,0.00}{##1}}}
\@namedef{PY@tok@kd}{\let\PY@bf=\textbf\def\PY@tc##1{\textcolor[rgb]{0.00,0.50,0.00}{##1}}}
\@namedef{PY@tok@kn}{\let\PY@bf=\textbf\def\PY@tc##1{\textcolor[rgb]{0.00,0.50,0.00}{##1}}}
\@namedef{PY@tok@kr}{\let\PY@bf=\textbf\def\PY@tc##1{\textcolor[rgb]{0.00,0.50,0.00}{##1}}}
\@namedef{PY@tok@bp}{\def\PY@tc##1{\textcolor[rgb]{0.00,0.50,0.00}{##1}}}
\@namedef{PY@tok@fm}{\def\PY@tc##1{\textcolor[rgb]{0.00,0.00,1.00}{##1}}}
\@namedef{PY@tok@vc}{\def\PY@tc##1{\textcolor[rgb]{0.10,0.09,0.49}{##1}}}
\@namedef{PY@tok@vg}{\def\PY@tc##1{\textcolor[rgb]{0.10,0.09,0.49}{##1}}}
\@namedef{PY@tok@vi}{\def\PY@tc##1{\textcolor[rgb]{0.10,0.09,0.49}{##1}}}
\@namedef{PY@tok@vm}{\def\PY@tc##1{\textcolor[rgb]{0.10,0.09,0.49}{##1}}}
\@namedef{PY@tok@sa}{\def\PY@tc##1{\textcolor[rgb]{0.73,0.13,0.13}{##1}}}
\@namedef{PY@tok@sb}{\def\PY@tc##1{\textcolor[rgb]{0.73,0.13,0.13}{##1}}}
\@namedef{PY@tok@sc}{\def\PY@tc##1{\textcolor[rgb]{0.73,0.13,0.13}{##1}}}
\@namedef{PY@tok@dl}{\def\PY@tc##1{\textcolor[rgb]{0.73,0.13,0.13}{##1}}}
\@namedef{PY@tok@s2}{\def\PY@tc##1{\textcolor[rgb]{0.73,0.13,0.13}{##1}}}
\@namedef{PY@tok@sh}{\def\PY@tc##1{\textcolor[rgb]{0.73,0.13,0.13}{##1}}}
\@namedef{PY@tok@s1}{\def\PY@tc##1{\textcolor[rgb]{0.73,0.13,0.13}{##1}}}
\@namedef{PY@tok@mb}{\def\PY@tc##1{\textcolor[rgb]{0.40,0.40,0.40}{##1}}}
\@namedef{PY@tok@mf}{\def\PY@tc##1{\textcolor[rgb]{0.40,0.40,0.40}{##1}}}
\@namedef{PY@tok@mh}{\def\PY@tc##1{\textcolor[rgb]{0.40,0.40,0.40}{##1}}}
\@namedef{PY@tok@mi}{\def\PY@tc##1{\textcolor[rgb]{0.40,0.40,0.40}{##1}}}
\@namedef{PY@tok@il}{\def\PY@tc##1{\textcolor[rgb]{0.40,0.40,0.40}{##1}}}
\@namedef{PY@tok@mo}{\def\PY@tc##1{\textcolor[rgb]{0.40,0.40,0.40}{##1}}}
\@namedef{PY@tok@ch}{\let\PY@it=\textit\def\PY@tc##1{\textcolor[rgb]{0.24,0.48,0.48}{##1}}}
\@namedef{PY@tok@cm}{\let\PY@it=\textit\def\PY@tc##1{\textcolor[rgb]{0.24,0.48,0.48}{##1}}}
\@namedef{PY@tok@cpf}{\let\PY@it=\textit\def\PY@tc##1{\textcolor[rgb]{0.24,0.48,0.48}{##1}}}
\@namedef{PY@tok@c1}{\let\PY@it=\textit\def\PY@tc##1{\textcolor[rgb]{0.24,0.48,0.48}{##1}}}
\@namedef{PY@tok@cs}{\let\PY@it=\textit\def\PY@tc##1{\textcolor[rgb]{0.24,0.48,0.48}{##1}}}

\def\PYZbs{\char`\\}
\def\PYZus{\char`\_}
\def\PYZob{\char`\{}
\def\PYZcb{\char`\}}
\def\PYZca{\char`\^}
\def\PYZam{\char`\&}
\def\PYZlt{\char`\<}
\def\PYZgt{\char`\>}
\def\PYZsh{\char`\#}
\def\PYZpc{\char`\%}
\def\PYZdl{\char`\$}
\def\PYZhy{\char`\-}
\def\PYZsq{\char`\'}
\def\PYZdq{\char`\"}
\def\PYZti{\char`\~}
% for compatibility with earlier versions
\def\PYZat{@}
\def\PYZlb{[}
\def\PYZrb{]}
\makeatother


    % For linebreaks inside Verbatim environment from package fancyvrb.
    \makeatletter
        \newbox\Wrappedcontinuationbox
        \newbox\Wrappedvisiblespacebox
        \newcommand*\Wrappedvisiblespace {\textcolor{red}{\textvisiblespace}}
        \newcommand*\Wrappedcontinuationsymbol {\textcolor{red}{\llap{\tiny$\m@th\hookrightarrow$}}}
        \newcommand*\Wrappedcontinuationindent {3ex }
        \newcommand*\Wrappedafterbreak {\kern\Wrappedcontinuationindent\copy\Wrappedcontinuationbox}
        % Take advantage of the already applied Pygments mark-up to insert
        % potential linebreaks for TeX processing.
        %        {, <, #, %, $, ' and ": go to next line.
        %        _, }, ^, &, >, - and ~: stay at end of broken line.
        % Use of \textquotesingle for straight quote.
        \newcommand*\Wrappedbreaksatspecials {%
            \def\PYGZus{\discretionary{\char`\_}{\Wrappedafterbreak}{\char`\_}}%
            \def\PYGZob{\discretionary{}{\Wrappedafterbreak\char`\{}{\char`\{}}%
            \def\PYGZcb{\discretionary{\char`\}}{\Wrappedafterbreak}{\char`\}}}%
            \def\PYGZca{\discretionary{\char`\^}{\Wrappedafterbreak}{\char`\^}}%
            \def\PYGZam{\discretionary{\char`\&}{\Wrappedafterbreak}{\char`\&}}%
            \def\PYGZlt{\discretionary{}{\Wrappedafterbreak\char`\<}{\char`\<}}%
            \def\PYGZgt{\discretionary{\char`\>}{\Wrappedafterbreak}{\char`\>}}%
            \def\PYGZsh{\discretionary{}{\Wrappedafterbreak\char`\#}{\char`\#}}%
            \def\PYGZpc{\discretionary{}{\Wrappedafterbreak\char`\%}{\char`\%}}%
            \def\PYGZdl{\discretionary{}{\Wrappedafterbreak\char`\$}{\char`\$}}%
            \def\PYGZhy{\discretionary{\char`\-}{\Wrappedafterbreak}{\char`\-}}%
            \def\PYGZsq{\discretionary{}{\Wrappedafterbreak\textquotesingle}{\textquotesingle}}%
            \def\PYGZdq{\discretionary{}{\Wrappedafterbreak\char`\"}{\char`\"}}%
            \def\PYGZti{\discretionary{\char`\~}{\Wrappedafterbreak}{\char`\~}}%
        }
        % Some characters . , ; ? ! / are not pygmentized.
        % This macro makes them "active" and they will insert potential linebreaks
        \newcommand*\Wrappedbreaksatpunct {%
            \lccode`\~`\.\lowercase{\def~}{\discretionary{\hbox{\char`\.}}{\Wrappedafterbreak}{\hbox{\char`\.}}}%
            \lccode`\~`\,\lowercase{\def~}{\discretionary{\hbox{\char`\,}}{\Wrappedafterbreak}{\hbox{\char`\,}}}%
            \lccode`\~`\;\lowercase{\def~}{\discretionary{\hbox{\char`\;}}{\Wrappedafterbreak}{\hbox{\char`\;}}}%
            \lccode`\~`\:\lowercase{\def~}{\discretionary{\hbox{\char`\:}}{\Wrappedafterbreak}{\hbox{\char`\:}}}%
            \lccode`\~`\?\lowercase{\def~}{\discretionary{\hbox{\char`\?}}{\Wrappedafterbreak}{\hbox{\char`\?}}}%
            \lccode`\~`\!\lowercase{\def~}{\discretionary{\hbox{\char`\!}}{\Wrappedafterbreak}{\hbox{\char`\!}}}%
            \lccode`\~`\/\lowercase{\def~}{\discretionary{\hbox{\char`\/}}{\Wrappedafterbreak}{\hbox{\char`\/}}}%
            \catcode`\.\active
            \catcode`\,\active
            \catcode`\;\active
            \catcode`\:\active
            \catcode`\?\active
            \catcode`\!\active
            \catcode`\/\active
            \lccode`\~`\~
        }
    \makeatother

    \let\OriginalVerbatim=\Verbatim
    \makeatletter
    \renewcommand{\Verbatim}[1][1]{%
        %\parskip\z@skip
        \sbox\Wrappedcontinuationbox {\Wrappedcontinuationsymbol}%
        \sbox\Wrappedvisiblespacebox {\FV@SetupFont\Wrappedvisiblespace}%
        \def\FancyVerbFormatLine ##1{\hsize\linewidth
            \vtop{\raggedright\hyphenpenalty\z@\exhyphenpenalty\z@
                \doublehyphendemerits\z@\finalhyphendemerits\z@
                \strut ##1\strut}%
        }%
        % If the linebreak is at a space, the latter will be displayed as visible
        % space at end of first line, and a continuation symbol starts next line.
        % Stretch/shrink are however usually zero for typewriter font.
        \def\FV@Space {%
            \nobreak\hskip\z@ plus\fontdimen3\font minus\fontdimen4\font
            \discretionary{\copy\Wrappedvisiblespacebox}{\Wrappedafterbreak}
            {\kern\fontdimen2\font}%
        }%

        % Allow breaks at special characters using \PYG... macros.
        \Wrappedbreaksatspecials
        % Breaks at punctuation characters . , ; ? ! and / need catcode=\active
        \OriginalVerbatim[#1,codes*=\Wrappedbreaksatpunct]%
    }
    \makeatother

    % Exact colors from NB
    \definecolor{incolor}{HTML}{303F9F}
    \definecolor{outcolor}{HTML}{D84315}
    \definecolor{cellborder}{HTML}{CFCFCF}
    \definecolor{cellbackground}{HTML}{F7F7F7}

    % prompt
    \makeatletter
    \newcommand{\boxspacing}{\kern\kvtcb@left@rule\kern\kvtcb@boxsep}
    \makeatother
    \newcommand{\prompt}[4]{
        {\ttfamily\llap{{\color{#2}[#3]:\hspace{3pt}#4}}\vspace{-\baselineskip}}
    }
    

    
    % Prevent overflowing lines due to hard-to-break entities
    \sloppy
    % Setup hyperref package
    \hypersetup{
      breaklinks=true,  % so long urls are correctly broken across lines
      colorlinks=true,
      urlcolor=urlcolor,
      linkcolor=linkcolor,
      citecolor=citecolor,
      }
    % Slightly bigger margins than the latex defaults
    
    \geometry{verbose,tmargin=1in,bmargin=1in,lmargin=1in,rmargin=1in}
    
    

\begin{document}
    
    \maketitle
    
    

    
    \section{Python
数据数学分析核心讲义}\label{python-ux6570ux636eux6570ux5b66ux5206ux6790ux6838ux5fc3ux8bb2ux4e49}

    \subsection{引言:从工程代码回归数据本质}\label{ux5f15ux8a00ux4eceux5de5ux7a0bux4ee3ux7801ux56deux5f52ux6570ux636eux672cux8d28}

    \subsubsection{为什么是 Python?}\label{ux4e3aux4ec0ux4e48ux662f-python}

    在传统的软件后端开发中, 关注点通常是 \textbf{状态管理、网络
I/O、并发安全与面向对象的解耦抽象}. 然而,
当进入数据分析与科学计算领域时, 思维模型需要发生根本性转变:

\begin{verbatim}
[传统后端工程思维]                     [现代数据分析/矩阵思维]
对象封装(OOP)                    ->   数据与操作分离(Data-Oriented Design)
逐个遍历(for item in list)       ->   批量向量化计算(SIMD / Tensor Processing)
复杂控制流(if-else / 状态机)      ->   布尔掩码与索引选择(Mask & Fancy Indexing)
显式多线程/锁竞争                 ->   底层 C/Fortran/BLAS 内存连续并行计算
\end{verbatim}

Python 之所以能够统治现代数据科学与 AI 领域,
并非因为其解释器运行效率高(相反, CPython 的解释器开销与 GIL
限制众所周知), 而是因为它的 \textbf{``胶水特性'' 与 C-API 扩展能力}:

\begin{enumerate}
\def\labelenumi{\arabic{enumi}.}
\tightlist
\item
  上层提供极具表现力的动态语法和交互式 REPL 环境;
\item
  下层无缝对接 OpenBLAS、MKL、CUDA 等底层硬件加速库;
\item
  形成了从数据摄取(Pandas/Arrow)、处理(NumPy/SciPy)、可视化(Matplotlib/Seaborn)到模型训练(PyTorch/Scikit-learn)的完整闭环.
\end{enumerate}

\begin{quote}
为什么 Python 明明什么都能算, 我们还需要 NumPy 和 Pandas?
\end{quote}

\begin{quote}
NumPy 和 Pandas 真正带来的价值, 是``多了一堆 API'',
还是``改变了数据和计算的抽象方式''?
\end{quote}

    \subsubsection{我们真正要学的不是
API}\label{ux6211ux4eecux771fux6b63ux8981ux5b66ux7684ux4e0dux662f-api}

    假设现在拿到 200 万条出租车订单.

我们需要分析:

\begin{Shaded}
\begin{Highlighting}[]
\NormalTok{1. 哪些订单不可信? }
\NormalTok{2. 什么时候最忙? }
\NormalTok{3. 什么时候流水最高? }
\NormalTok{4. 哪些区域最热门? }
\NormalTok{5. 哪些订单看起来异常? }
\NormalTok{6. 为什么同一个计算, 有人的代码跑几十秒, 有人的代码只跑几秒? }
\NormalTok{7. 怎么把这些分析变成公司里以后还能复用的工具? }
\end{Highlighting}
\end{Shaded}

当然, 可以全部使用 Python:

\begin{Shaded}
\begin{Highlighting}[]
\ControlFlowTok{for}\NormalTok{ row }\KeywordTok{in}\NormalTok{ rows:}
\NormalTok{    ...}
\end{Highlighting}
\end{Shaded}

但问题并不在于 Python \textbf{能不能做}.

真正的问题是:

\begin{quote}
\textbf{Python 原生的数据结构和逐对象执行模型,
并不是专门针对百万级同构数值计算和二维表格分析设计的.}
\end{quote}

但是 Numpy 和 Pandas 是 Python 中专门为数据分析设计的

\begin{enumerate}
\def\labelenumi{\arabic{enumi}.}
\tightlist
\item
  NumPy 的核心是同构 \(N\) 维数组 \texttt{ndarray}; 官方文档将其定义为
  NumPy 的核心多维数组结构,
  并围绕数组提供高效的数学、逻辑、选择、排序、线性代数等运算.
\item
  Pandas 则在数组之上增加了
  \textbf{标签、索引、缺失值、异构列、自动对齐、分组、连接和时间序列}
  等更贴近表格数据分析的抽象; \texttt{Index} 本身就是 Pandas
  用于索引和对齐的轴标签对象.
\end{enumerate}

{\def\LTcaptype{none} % do not increment counter
\begin{longtable}[]{@{}
  >{\centering\arraybackslash}p{(\linewidth - 6\tabcolsep) * \real{0.1919}}
  >{\centering\arraybackslash}p{(\linewidth - 6\tabcolsep) * \real{0.4444}}
  >{\centering\arraybackslash}p{(\linewidth - 6\tabcolsep) * \real{0.2525}}
  >{\centering\arraybackslash}p{(\linewidth - 6\tabcolsep) * \real{0.1111}}@{}}
\toprule\noalign{}
\begin{minipage}[b]{\linewidth}\centering
从哪里来
\end{minipage} & \begin{minipage}[b]{\linewidth}\centering
核心技术
\end{minipage} & \begin{minipage}[b]{\linewidth}\centering
关键机制
\end{minipage} & \begin{minipage}[b]{\linewidth}\centering
最终走向
\end{minipage} \\
\midrule\noalign{}
\endhead
\bottomrule\noalign{}
\endlastfoot
\textbf{Python 原生对象} & \texttt{list} / \texttt{dict} / Object &
灵活，但对象开销大、逐元素处理 & ↓ \\
\textbf{NumPy} & \texttt{ndarray} &
\texttt{dtype}、\texttt{shape}、\texttt{strides} &
\textbf{高效数值计算} \\
\textbf{NumPy 计算模型} & UFunc / Vectorization / Broadcasting &
批量计算、底层循环、形状扩展 & \textbf{性能提升} \\
\textbf{Pandas} & Series / DataFrame & Index、Alignment &
\textbf{结构化数据分析} \\
\textbf{Pandas 数据操作} & \texttt{groupby} / \texttt{merge} /
\texttt{resample} / \texttt{rolling} & 聚合、关联、时间序列、窗口计算 &
\textbf{复杂数据处理} \\
\textbf{业务分析} & 数据清洗 → 分析 → 建模 & 将工具应用于真实数据 &
\textbf{业务结论} \\
\textbf{性能工程} & 内存、向量化、拷贝、计算瓶颈 & 针对大规模数据优化 &
\textbf{更快、更省内存} \\
\textbf{领域工具 / Pipeline} & 专用工具、Accessor、Pipeline &
封装通用能力 & \textbf{工程化应用} \\
\end{longtable}
}

一句话概括: \textgreater{} \textbf{NumPy 解决``如何高效计算数组'';
Pandas 解决``如何带着业务语义高效组织、查询、组合和分析表格数据''.}

    \subsubsection{出租车数据集简介}\label{ux51faux79dfux8f66ux6570ux636eux96c6ux7b80ux4ecb}

    NYC TLC Yellow Taxi Trip Records 是纽约市出租车与豪华轿车委员会(NYC Taxi
\& Limousine Commission, TLC)公开发布的纽约市黄色出租车行程记录数据.

{\def\LTcaptype{none} % do not increment counter
\begin{longtable}[]{@{}
  >{\centering\arraybackslash}p{(\linewidth - 6\tabcolsep) * \real{0.1139}}
  >{\centering\arraybackslash}p{(\linewidth - 6\tabcolsep) * \real{0.2911}}
  >{\centering\arraybackslash}p{(\linewidth - 6\tabcolsep) * \real{0.3291}}
  >{\centering\arraybackslash}p{(\linewidth - 6\tabcolsep) * \real{0.2658}}@{}}
\toprule\noalign{}
\begin{minipage}[b]{\linewidth}\centering
类别
\end{minipage} & \begin{minipage}[b]{\linewidth}\centering
字段
\end{minipage} & \begin{minipage}[b]{\linewidth}\centering
含义
\end{minipage} & \begin{minipage}[b]{\linewidth}\centering
典型值/单位
\end{minipage} \\
\midrule\noalign{}
\endhead
\bottomrule\noalign{}
\endlastfoot
\textbf{记录信息} & \texttt{VendorID} &
提供出租车计价/数据记录服务的技术供应商编号 & \texttt{1}, \texttt{2},
\texttt{6}, \texttt{7} \\
\textbf{时间} & \texttt{tpep\_pickup\_datetime} &
计价器开始计费的时间，即上车时间 & \texttt{2023-01-01\ 00:15:00} \\
& \texttt{tpep\_dropoff\_datetime} & 计价器停止计费的时间，即下车时间 &
\texttt{2023-01-01\ 00:25:00} \\
\textbf{乘客/行程} & \texttt{passenger\_count} & 乘客人数 & \texttt{1},
\texttt{2}, \texttt{3}\ldots{} \\
& \texttt{trip\_distance} & 计价器记录的行程距离 & 英里 mile \\
\textbf{费率} & \texttt{RatecodeID} & 本次行程最终使用的费率类型 &
\texttt{1}, \texttt{2}, \texttt{3}\ldots{} \\
\textbf{数据传输} & \texttt{store\_and\_fwd\_flag} &
行程数据是否曾暂存在车载设备中后再上传 & \texttt{Y} / \texttt{N} \\
\textbf{空间位置} & \texttt{PULocationID} & 上车所在的 Taxi Zone & Zone
ID \\
& \texttt{DOLocationID} & 下车所在的 Taxi Zone & Zone ID \\
\textbf{支付} & \texttt{payment\_type} & 支付方式 & 信用卡、现金等 \\
\textbf{费用} & \texttt{fare\_amount} & 计价器计算出的基础车费 & 美元 \\
& \texttt{extra} & 额外费用/附加费 & 美元 \\
& \texttt{mta\_tax} & MTA 税费 & 美元 \\
& \texttt{tip\_amount} & 小费金额 & 美元 \\
& \texttt{tolls\_amount} & 行程产生的过路费 & 美元 \\
& \texttt{improvement\_surcharge} & Taxi Improvement Surcharge & 美元 \\
& \texttt{total\_amount} & 向乘客收取的总金额 & 美元 \\
& \texttt{congestion\_surcharge} & 拥堵附加费 & 美元 \\
& \texttt{airport\_fee} & JFK / LaGuardia 机场相关费用 & 美元 \\
& \texttt{cbd\_congestion\_fee} & CBD 拥堵费，2025 年起出现 & 美元 \\
\end{longtable}
}

    \paragraph{数据的格式}\label{ux6570ux636eux7684ux683cux5f0f}

    {\def\LTcaptype{none} % do not increment counter
\begin{longtable}[]{@{}
  >{\raggedright\arraybackslash}p{(\linewidth - 4\tabcolsep) * \real{0.3333}}
  >{\raggedright\arraybackslash}p{(\linewidth - 4\tabcolsep) * \real{0.3333}}
  >{\raggedright\arraybackslash}p{(\linewidth - 4\tabcolsep) * \real{0.3333}}@{}}
\toprule\noalign{}
\begin{minipage}[b]{\linewidth}\raggedright
特性维度
\end{minipage} & \begin{minipage}[b]{\linewidth}\raggedright
CSV
\end{minipage} & \begin{minipage}[b]{\linewidth}\raggedright
Parquet
\end{minipage} \\
\midrule\noalign{}
\endhead
\bottomrule\noalign{}
\endlastfoot
\textbf{存储结构} & 行式存储(Row-oriented，纯文本) &
列式存储(Columnar，二进制) \\
\textbf{文件体积} & 较大(无内置压缩，字符冗余高) & 极小(内置 Snappy/ZSTD
等高效压缩与字典编码) \\
\textbf{读取性能} & 较慢(需逐行扫描、切分字符串并推断类型) &
极快(支持列裁剪与分区下推) \\
\textbf{Schema \& 类型} & 弱类型(无元数据，需解析器自行推断) &
强类型(内嵌 Schema、数据类型及统计信息) \\
\textbf{可读性} & 极高(文本编辑器、Excel 直接打开修改) &
无法直接阅读(需专用解析引擎) \\
\textbf{I/O 优化} & 必须全表/全行扫描 & 仅读取查询所需的列和数据块(Row
Groups) \\
\end{longtable}
}

    \subsubsection{贯穿本节的七个问题}\label{ux8d2fux7a7fux672cux8282ux7684ux4e03ux4e2aux95eeux9898}

    {\def\LTcaptype{none} % do not increment counter
\begin{longtable}[]{@{}
  >{\centering\arraybackslash}p{(\linewidth - 4\tabcolsep) * \real{0.3333}}
  >{\centering\arraybackslash}p{(\linewidth - 4\tabcolsep) * \real{0.3333}}
  >{\centering\arraybackslash}p{(\linewidth - 4\tabcolsep) * \real{0.3333}}@{}}
\toprule\noalign{}
\begin{minipage}[b]{\linewidth}\centering
问题
\end{minipage} & \begin{minipage}[b]{\linewidth}\centering
业务问题
\end{minipage} & \begin{minipage}[b]{\linewidth}\centering
逐步引出的技术
\end{minipage} \\
\midrule\noalign{}
\endhead
\bottomrule\noalign{}
\endlastfoot
\textbf{Q1} & 这 200 万条数据可靠吗? &
\texttt{dtype}、缺失值、mask、\texttt{loc}、\texttt{query} \\
\textbf{Q2} & 一天中什么时候最忙? &
\texttt{datetime}、\texttt{groupby}、\texttt{agg} \\
\textbf{Q3} & 什么时间段的流水和单位运营时间收入最高? &
向量化、\texttt{assign}、\texttt{groupby}、派生指标 \\
\textbf{Q4} & 哪些上下车区域最热门? &
\texttt{value\_counts}、\texttt{merge}、连接校验 \\
\textbf{Q5} & 哪些订单和时间窗口可能异常? &
\texttt{where}、\texttt{select}、\texttt{quantile}、\texttt{transform}、\texttt{resample}、\texttt{rolling} \\
\textbf{Q6} & 为什么同一个分析可以相差一个数量级甚至更多? &
\texttt{Python\ loop}、\texttt{apply}、\texttt{NumPy}、\texttt{Numba}、\texttt{eval/query}、内存优化 \\
\textbf{Q7} & 如何把分析脚本变成公司可以重复使用的工具? &
\texttt{pipe}、纯函数、\texttt{Accessor}、领域 API \\
\end{longtable}
}

    \subsubsection{NumPy / Pandas
为什么存在}\label{numpy-pandas-ux4e3aux4ec0ux4e48ux5b58ux5728}

    \paragraph{问题}\label{ux95eeux9898}

    假设我们有 200 万个金额数据:

\begin{Shaded}
\begin{Highlighting}[]
\NormalTok{prices }\OperatorTok{=}\NormalTok{ [...]}
\end{Highlighting}
\end{Shaded}

需要:

\begin{Shaded}
\begin{Highlighting}[]
\NormalTok{全部上涨 5\%}
\NormalTok{再增加 3 元服务费}
\NormalTok{过滤异常价格}
\NormalTok{求平均值}
\end{Highlighting}
\end{Shaded}

    \paragraph{直接使用 Python List
实现}\label{ux76f4ux63a5ux4f7fux7528-python-list-ux5b9eux73b0}

    \begin{Shaded}
\begin{Highlighting}[]
\CommentTok{\# 直接使用 List 存储与计算}
\NormalTok{result }\OperatorTok{=}\NormalTok{ []}

\ControlFlowTok{for}\NormalTok{ price }\KeywordTok{in}\NormalTok{ prices:}
\NormalTok{    result.append(price }\OperatorTok{*} \FloatTok{1.05} \OperatorTok{+} \DecValTok{3}\NormalTok{)}
\end{Highlighting}
\end{Shaded}

问题是我们把 \textbf{200 万次循环控制} 交给了 Python 层. 具体来说,
Python 原生 \texttt{List} 在这种大批量数值计算场景下存在两大明显短板:

\textbf{1. 存储方面}

Python \texttt{List} 并不是一块连续的数值缓冲区, 而是一个
\textbf{对象引用数组}. 它存储的每一个元素都是一个独立的 Python
对象(这里是 \texttt{float}), 每个对象都带有引用计数、类型指针等元信息.
以 200 万个金额为例:

\begin{Shaded}
\begin{Highlighting}[]
\NormalTok{prices }\OperatorTok{=}\NormalTok{ [}\FloatTok{42.5}\NormalTok{, }\FloatTok{17.3}\NormalTok{, }\FloatTok{58.9}\NormalTok{, ...]   }\CommentTok{\# 200 万个 float 对象}
\end{Highlighting}
\end{Shaded}

\begin{itemize}
\tightlist
\item
  \texttt{List} 本身只存 200 万个\textbf{指针}, 每个指针 8 字节(64
  位系统);
\item
  每一个 \texttt{float} 对象本身还要额外占用约 24 字节(CPython 对象头);
\item
  于是 \texttt{List} 存储 200 万个浮点数的实际内存, 大约是相同规模
  \texttt{ndarray} 的 \textbf{3\textasciitilde5 倍}.
\end{itemize}

\begin{Shaded}
\begin{Highlighting}[]
\NormalTok{List:}
\NormalTok{  指针数组(8B × 200万) + 每个 float 对象(24B × 200万)}
\NormalTok{  ≈ 16 MB + 48 MB = 64 MB 左右}

\NormalTok{ndarray(float64):}
\NormalTok{  连续缓冲区(8B × 200万)}
\NormalTok{  ≈ 16 MB}
\end{Highlighting}
\end{Shaded}

内存翻倍, 还意味着缓存命中率更低、CPU 需要跨越更多随机内存地址.

\textbf{2. 计算方面}

\begin{Shaded}
\begin{Highlighting}[]
\NormalTok{result }\OperatorTok{=}\NormalTok{ []}
\ControlFlowTok{for}\NormalTok{ price }\KeywordTok{in}\NormalTok{ prices:}
\NormalTok{    result.append(price }\OperatorTok{*} \FloatTok{1.05} \OperatorTok{+} \DecValTok{3}\NormalTok{)}
\end{Highlighting}
\end{Shaded}

这段循环的执行, 每一步都发生在 \textbf{Python 解释器层},
而不是底层机器码层:

\begin{itemize}
\tightlist
\item
  每次 \texttt{for} 迭代都要经历字节码分发(bytecode dispatch);
\item
  每次读取 \texttt{price} 都要做一次 \textbf{拆箱}(unboxing, 从 Python
  对象还原成 C double);
\item
  每次 \texttt{price\ *\ 1.05\ +\ 3} 结果要重新
  \textbf{装箱}(boxing)成新的 \texttt{float} 对象;
\item
  每次 \texttt{append}
  还要调用方法、检查容量、可能触发列表扩容与内存重分配.
\end{itemize}

当循环次数上升到 200 万, 这些\textbf{解释器级开销被放大 200 万倍},
成为无可忽视的性能瓶颈. 这也是为什么同样的计算, 纯 Python 循环往往比
NumPy 向量化写法慢一到两个数量级------\textbf{真正慢的不是''算'',
而是''每次都要让解释器来驱动''}.

    \paragraph{使用 Numpy 和 Pandas
实现}\label{ux4f7fux7528-numpy-ux548c-pandas-ux5b9eux73b0}

    NumPy:

\begin{Shaded}
\begin{Highlighting}[]
\NormalTok{result }\OperatorTok{=}\NormalTok{ prices }\OperatorTok{*} \FloatTok{1.05} \OperatorTok{+} \DecValTok{3}
\end{Highlighting}
\end{Shaded}

Pandas:

\begin{Shaded}
\begin{Highlighting}[]
\NormalTok{result }\OperatorTok{=}\NormalTok{ (}
\NormalTok{    df}
\NormalTok{    .query(}\StringTok{"fare\_amount \textgreater{} 0"}\NormalTok{)}
\NormalTok{    .groupby(}\StringTok{"pickup\_hour"}\NormalTok{)}
\NormalTok{    .agg(avg\_fare}\OperatorTok{=}\NormalTok{(}\StringTok{"fare\_amount"}\NormalTok{, }\StringTok{"mean"}\NormalTok{))}
\NormalTok{)}
\end{Highlighting}
\end{Shaded}

    \subsubsection{Numpy 和 Pandas
的特点}\label{numpy-ux548c-pandas-ux7684ux7279ux70b9}

    \textbf{NumPy 解决的痛点(纯数值与张量计算):} 在科学计算场景下, 原生
Python 的 \texttt{list}
存储机制会导致严重的内存碎片化(频繁指针解引用)和缓存未命中(Cache Miss),
同时动态类型解释器导致循环极慢. NumPy 的诞生终结了这五大痛点:
解决内存空间碎片化、消除解释器动态类型开销、填补多维张量代数的空白、终结切片的高昂拷贝代价,
并打通底层 C/BLAS 和 SIMD 硬件加速生态.

\textbf{Pandas 解决的痛点(异构业务数据与关系代数):}
真实世界的业务数据不是纯数学矩阵,
而是包含字符串、时间戳、布尔值的异构表格, 常存在缺失值和错位情况.Pandas
解决了: 1. 如何容纳异构数据类型(字符串、数值混合). 2.
在连接表或对齐时如何依据业务``键(Label)''而非纯``位置(Index)''进行运算.
3. 如何像 SQL 一样使用 GroupBy、Join 等关系代数算子. 4.
如何优雅且不中断程序地处理 NaN 缺失数据.

    \paragraph{核心设计哲学}\label{ux6838ux5fc3ux8bbeux8ba1ux54f2ux5b66}

    \subparagraph{\texorpdfstring{NumPy
的灵魂核心:\texttt{ndarray}(N-dimensional
array)}{NumPy 的灵魂核心:ndarray(N-dimensional array)}}\label{numpy-ux7684ux7075ux9b42ux6838ux5fc3ndarrayn-dimensional-array}

    \texttt{ndarray}(N-Dimensional Array)
的本质是``同构多维连续内存缓冲区的视图抽象(N-dimensional Array /
Tensor)'', 矩阵(Matrix)是它在 2 维空间下的一个特例.
它的设计哲学是:\textbf{在动态语言中实现静态语言级别的内存布局.}

\begin{itemize}
\tightlist
\item
  \textbf{元数据解耦与视图(View):} \texttt{ndarray} 将数据的
  \texttt{Header}(包含 shape、dtype、strides)与底层的
  \texttt{Data\ Buffer}(连续物理内存)彻底分离. 这使得切片、转置、Reshape
  都只需要极小开销修改 Header 且不需要复制物理内存.
\item
  \textbf{齐次强类型与 SIMD:} 数据块内类型必须一致, 使得 CPU
  能够跨步寻址(索引 × 大小), 并直接投喂给硬件向量化指令.
\item
  \textbf{面向数组编程:} 把底层长循环隐式推给 C 语言层面,
  上层仅做``张量形态''的声明式表达.
\end{itemize}

    \subparagraph{\texorpdfstring{Pandas 的根本基石:\texttt{Series} 和
\texttt{DataFrame}}{Pandas 的根本基石:Series 和 DataFrame}}\label{pandas-ux7684ux6839ux672cux57faux77f3series-ux548c-dataframe}

    如果 NumPy 是``同质数学张量'', Pandas 则是``带标签的异构列式数据框''. *
\textbf{列式存储架构:} DataFrame 本质是多个一维 \texttt{Series}
的字典集合. 在底层, 同类型的列会被组织为内存连续的块.
列式存储使得``求整列均值''等分析操作在物理上读取连续内存, 极其高效. *
\textbf{标签驱动(Label-Based Alignment):} 与 NumPy
严格的矩阵位置形状对齐不同, Series 和 DataFrame
的核心哲学在于``业务主键(Index)绑定''. 两张表相加时, Pandas
会自动依据相同的 Index 标签进行匹配(如果没有则视为 NaN),
哪怕它们数据行数或顺序完全不同. 这规避了业务分析中最致命的错位问题.

    \subsection{课程环境与初始化}\label{ux8bfeux7a0bux73afux5883ux4e0eux521dux59cbux5316}

    \subsubsection{Numpy 与 Pandas
的版本}\label{numpy-ux4e0e-pandas-ux7684ux7248ux672c}

    本课程 Numpy 与 Pandas 需要安装与配套的版本如下:

{\def\LTcaptype{none} % do not increment counter
\begin{longtable}[]{@{}cc@{}}
\toprule\noalign{}
Package & Version \\
\midrule\noalign{}
\endhead
\bottomrule\noalign{}
\endlastfoot
matplotlib & 3.9.4 \\
numba & 0.60.0 \\
numpy & 2.0.2 \\
pandas & 2.3.3 \\
pyarrow & 21.0.0 \\
jupyterlab & 4.5.10 \\
\end{longtable}
}

    \subsubsection{本课程统一 Benchmark
工具}\label{ux672cux8bfeux7a0bux7edfux4e00-benchmark-ux5de5ux5177}

    \subsection{NumPy:真正理解
ndarray}\label{numpyux771fux6b63ux7406ux89e3-ndarray}

    通过这部分希望大家可以理解下面的问题:

\begin{Shaded}
\begin{Highlighting}[]
\NormalTok{ndarray ≠ 更快的 list}

\NormalTok{ndarray =}
\NormalTok{    data buffer}
\NormalTok{    + dtype}
\NormalTok{    + shape}
\NormalTok{    + strides}
\end{Highlighting}
\end{Shaded}

能够解释:

\begin{itemize}
\tightlist
\item
  \texttt{dtype}
\item
  \texttt{shape}
\item
  \texttt{ndim}
\item
  \texttt{size}
\item
  \texttt{itemsize}
\item
  \texttt{nbytes}
\item
  \texttt{strides}
\item
  \texttt{flags}
\item
  \texttt{view}
\item
  \texttt{copy}
\item
  \texttt{reshape}
\item
  \texttt{transpose}
\item
  \texttt{axis}
\end{itemize}

    \subsubsection{ndarray 是什么}\label{ndarray-ux662fux4ec0ux4e48}

    ndarray 是 Numpy 中的一个基础类型, 它底层的简化定义如下:

\begin{Shaded}
\begin{Highlighting}[]
\KeywordTok{typedef} \KeywordTok{struct} \OperatorTok{\{}
\NormalTok{    PyObject\_HEAD}

    \CommentTok{/* Number of dimensions */}
    \DataTypeTok{int}\NormalTok{ nd}\OperatorTok{;}
    \CommentTok{/* Shape of the array */}
\NormalTok{    npy\_intp }\OperatorTok{*}\NormalTok{dimensions}\OperatorTok{;}
    \CommentTok{/* Strides of the array */}
\NormalTok{    npy\_intp }\OperatorTok{*}\NormalTok{strides}\OperatorTok{;}
    \CommentTok{/* Pointer to the actual data */}
    \DataTypeTok{void} \OperatorTok{*}\NormalTok{data}\OperatorTok{;}
    \CommentTok{/* Data type descriptor */}
\NormalTok{    PyArray\_Descr }\OperatorTok{*}\NormalTok{descr}\OperatorTok{;}
    \CommentTok{/* Base object, used for views */}
\NormalTok{    PyObject }\OperatorTok{*}\NormalTok{base}\OperatorTok{;}
    \CommentTok{/* Flags */}
    \DataTypeTok{int}\NormalTok{ flags}\OperatorTok{;}
    \CommentTok{/* Weak references */}
\NormalTok{    PyObject }\OperatorTok{*}\NormalTok{weakreflist}\OperatorTok{;}
    \CommentTok{/* ... other internal fields ... */}

\OperatorTok{\}}\NormalTok{ PyArrayObject}\OperatorTok{;}
\end{Highlighting}
\end{Shaded}

NumPy 官方 \texttt{ndarray} 文档把
\texttt{shape}、\texttt{strides}、\texttt{ndim}、\texttt{data}、\texttt{itemsize}、\texttt{nbytes}、\texttt{base}、\texttt{dtype}
等直接列为数组(ndarray)的核心属性; \texttt{strides}
表示沿每个维度前进一步需要跨过的字节数.

{\def\LTcaptype{none} % do not increment counter
\begin{longtable}[]{@{}
  >{\centering\arraybackslash}p{(\linewidth - 4\tabcolsep) * \real{0.2632}}
  >{\centering\arraybackslash}p{(\linewidth - 4\tabcolsep) * \real{0.0877}}
  >{\centering\arraybackslash}p{(\linewidth - 4\tabcolsep) * \real{0.6491}}@{}}
\toprule\noalign{}
\begin{minipage}[b]{\linewidth}\centering
\texttt{ndarray} 核心组成
\end{minipage} & \begin{minipage}[b]{\linewidth}\centering
对应概念
\end{minipage} & \begin{minipage}[b]{\linewidth}\centering
含义
\end{minipage} \\
\midrule\noalign{}
\endhead
\bottomrule\noalign{}
\endlastfoot
\textbf{data buffer} & 数据缓冲区 & \textbf{真正存储数据的内存区域} \\
\textbf{dtype} & 数据类型 & 描述\textbf{一个元素应该怎样解释}，例如
\texttt{int64}、\texttt{float32} \\
\textbf{shape} & 形状 & 描述数组的\textbf{逻辑维度}，例如
\texttt{(3,\ 4)} 表示 3 行 4 列 \\
\textbf{strides} & 步长 & 描述\textbf{逻辑坐标如何映射到实际内存地址} \\
\end{longtable}
}

所以 \texttt{ndarray} 本质上可以理解为: 一块数据内存(Data Buffer) +
一组描述这块内存如何被解释和访问的元数据(dtype、shape、strides).

\begin{Shaded}
\begin{Highlighting}[]
\NormalTok{Python list}

\NormalTok{┌─────┐}
\NormalTok{│ ptr │──→ Python object}
\NormalTok{├─────┤}
\NormalTok{│ ptr │──→ Python object}
\NormalTok{├─────┤}
\NormalTok{│ ptr │──→ Python object}
\NormalTok{└─────┘}

\NormalTok{NumPy ndarray}

\NormalTok{metadata:}
\NormalTok{shape   = (4,)}
\NormalTok{dtype   = float64}
\NormalTok{strides = (8,)}

\NormalTok{data buffer:}
\NormalTok{┌────────┬────────┬────────┬────────┐}
\NormalTok{│ 1.2    │ 3.4    │ 5.6    │ 7.8    │}
\NormalTok{└────────┴────────┴────────┴────────┘}
\end{Highlighting}
\end{Shaded}

    \paragraph{新建一个 ndarray}\label{ux65b0ux5efaux4e00ux4e2a-ndarray}

我们创建一个最基本的 ndarray 如下:

    \begin{tcolorbox}[breakable, size=fbox, boxrule=1pt, pad at break*=1mm,colback=cellbackground, colframe=cellborder]
\prompt{In}{incolor}{7}{\boxspacing}
\begin{Verbatim}[commandchars=\\\{\}]
\PY{n}{a} \PY{o}{=} \PY{n}{np}\PY{o}{.}\PY{n}{array}\PY{p}{(}\PY{p}{[}\PY{l+m+mi}{1}\PY{p}{,} \PY{l+m+mi}{2}\PY{p}{,} \PY{l+m+mi}{3}\PY{p}{,} \PY{l+m+mi}{4}\PY{p}{]}\PY{p}{,} \PY{n}{dtype}\PY{o}{=}\PY{n}{np}\PY{o}{.}\PY{n}{int32}\PY{p}{)}

\PY{n+nb}{print}\PY{p}{(}\PY{l+s+s2}{\PYZdq{}}\PY{l+s+s2}{array   :}\PY{l+s+s2}{\PYZdq{}}\PY{p}{,} \PY{n}{a}\PY{p}{)}
\PY{n+nb}{print}\PY{p}{(}\PY{l+s+s2}{\PYZdq{}}\PY{l+s+s2}{dtype   :}\PY{l+s+s2}{\PYZdq{}}\PY{p}{,} \PY{n}{a}\PY{o}{.}\PY{n}{dtype}\PY{p}{)}
\PY{n+nb}{print}\PY{p}{(}\PY{l+s+s2}{\PYZdq{}}\PY{l+s+s2}{shape   :}\PY{l+s+s2}{\PYZdq{}}\PY{p}{,} \PY{n}{a}\PY{o}{.}\PY{n}{shape}\PY{p}{)}
\PY{n+nb}{print}\PY{p}{(}\PY{l+s+s2}{\PYZdq{}}\PY{l+s+s2}{ndim    :}\PY{l+s+s2}{\PYZdq{}}\PY{p}{,} \PY{n}{a}\PY{o}{.}\PY{n}{ndim}\PY{p}{)}
\PY{n+nb}{print}\PY{p}{(}\PY{l+s+s2}{\PYZdq{}}\PY{l+s+s2}{size    :}\PY{l+s+s2}{\PYZdq{}}\PY{p}{,} \PY{n}{a}\PY{o}{.}\PY{n}{size}\PY{p}{)}
\PY{n+nb}{print}\PY{p}{(}\PY{l+s+s2}{\PYZdq{}}\PY{l+s+s2}{itemsize:}\PY{l+s+s2}{\PYZdq{}}\PY{p}{,} \PY{n}{a}\PY{o}{.}\PY{n}{itemsize}\PY{p}{)}
\PY{n+nb}{print}\PY{p}{(}\PY{l+s+s2}{\PYZdq{}}\PY{l+s+s2}{nbytes  :}\PY{l+s+s2}{\PYZdq{}}\PY{p}{,} \PY{n}{a}\PY{o}{.}\PY{n}{nbytes}\PY{p}{)}
\PY{n+nb}{print}\PY{p}{(}\PY{l+s+s2}{\PYZdq{}}\PY{l+s+s2}{strides :}\PY{l+s+s2}{\PYZdq{}}\PY{p}{,} \PY{n}{a}\PY{o}{.}\PY{n}{strides}\PY{p}{)}

\PY{c+c1}{\PYZsh{} Numpy 内置了很多初始化函数: }
\PY{n}{m} \PY{o}{=} \PY{n}{np}\PY{o}{.}\PY{n}{arange}\PY{p}{(}\PY{l+m+mi}{12}\PY{p}{,} \PY{n}{dtype}\PY{o}{=}\PY{n}{np}\PY{o}{.}\PY{n}{int32}\PY{p}{)}\PY{o}{.}\PY{n}{reshape}\PY{p}{(}\PY{l+m+mi}{3}\PY{p}{,} \PY{l+m+mi}{4}\PY{p}{)}

\PY{n+nb}{print}\PY{p}{(}\PY{l+s+s2}{\PYZdq{}}\PY{l+s+s2}{shape   :}\PY{l+s+s2}{\PYZdq{}}\PY{p}{,} \PY{n}{m}\PY{o}{.}\PY{n}{shape}\PY{p}{)}
\PY{n+nb}{print}\PY{p}{(}\PY{l+s+s2}{\PYZdq{}}\PY{l+s+s2}{dtype   :}\PY{l+s+s2}{\PYZdq{}}\PY{p}{,} \PY{n}{m}\PY{o}{.}\PY{n}{dtype}\PY{p}{)}
\PY{n+nb}{print}\PY{p}{(}\PY{l+s+s2}{\PYZdq{}}\PY{l+s+s2}{strides :}\PY{l+s+s2}{\PYZdq{}}\PY{p}{,} \PY{n}{m}\PY{o}{.}\PY{n}{strides}\PY{p}{)}
\PY{n+nb}{print}\PY{p}{(}\PY{l+s+s2}{\PYZdq{}}\PY{l+s+s2}{flags:}\PY{l+s+s2}{\PYZdq{}}\PY{p}{)}
\PY{n+nb}{print}\PY{p}{(}\PY{n}{m}\PY{o}{.}\PY{n}{flags}\PY{p}{)}
\end{Verbatim}
\end{tcolorbox}

    \begin{Verbatim}[commandchars=\\\{\}]
array   : [1 2 3 4]
dtype   : int32
shape   : (4,)
ndim    : 1
size    : 4
itemsize: 4
nbytes  : 16
strides : (4,)
shape   : (3, 4)
dtype   : int32
strides : (16, 4)
flags:
  C\_CONTIGUOUS : True
  F\_CONTIGUOUS : False
  OWNDATA : False
  WRITEABLE : True
  ALIGNED : True
  WRITEBACKIFCOPY : False

    \end{Verbatim}

    对于上面的 C-order \texttt{int32} 数组, 一个元素 4 bytes:

\begin{Shaded}
\begin{Highlighting}[]
\NormalTok{shape = (3, 4)}

\NormalTok{内存:}

\NormalTok{0 1 2 3 | 4 5 6 7 | 8 9 10 11}

\NormalTok{列方向移动一步:4 bytes}
\NormalTok{行方向移动一步:4 × 4 = 16 bytes}

\NormalTok{strides = (16, 4)}
\end{Highlighting}
\end{Shaded}

常见创建函数:

    \begin{tcolorbox}[breakable, size=fbox, boxrule=1pt, pad at break*=1mm,colback=cellbackground, colframe=cellborder]
\prompt{In}{incolor}{8}{\boxspacing}
\begin{Verbatim}[commandchars=\\\{\}]
\PY{n}{examples} \PY{o}{=} \PY{p}{\PYZob{}}
    \PY{l+s+s2}{\PYZdq{}}\PY{l+s+s2}{array}\PY{l+s+s2}{\PYZdq{}}\PY{p}{:} \PY{n}{np}\PY{o}{.}\PY{n}{array}\PY{p}{(}\PY{p}{[}\PY{l+m+mi}{1}\PY{p}{,} \PY{l+m+mi}{2}\PY{p}{,} \PY{l+m+mi}{3}\PY{p}{]}\PY{p}{)}\PY{p}{,}
    \PY{l+s+s2}{\PYZdq{}}\PY{l+s+s2}{arange}\PY{l+s+s2}{\PYZdq{}}\PY{p}{:} \PY{n}{np}\PY{o}{.}\PY{n}{arange}\PY{p}{(}\PY{l+m+mi}{0}\PY{p}{,} \PY{l+m+mi}{10}\PY{p}{,} \PY{l+m+mi}{2}\PY{p}{)}\PY{p}{,}
    \PY{l+s+s2}{\PYZdq{}}\PY{l+s+s2}{linspace}\PY{l+s+s2}{\PYZdq{}}\PY{p}{:} \PY{n}{np}\PY{o}{.}\PY{n}{linspace}\PY{p}{(}\PY{l+m+mi}{0}\PY{p}{,} \PY{l+m+mi}{1}\PY{p}{,} \PY{l+m+mi}{5}\PY{p}{)}\PY{p}{,}
    \PY{l+s+s2}{\PYZdq{}}\PY{l+s+s2}{zeros}\PY{l+s+s2}{\PYZdq{}}\PY{p}{:} \PY{n}{np}\PY{o}{.}\PY{n}{zeros}\PY{p}{(}\PY{p}{(}\PY{l+m+mi}{2}\PY{p}{,} \PY{l+m+mi}{3}\PY{p}{)}\PY{p}{)}\PY{p}{,}
    \PY{l+s+s2}{\PYZdq{}}\PY{l+s+s2}{ones}\PY{l+s+s2}{\PYZdq{}}\PY{p}{:} \PY{n}{np}\PY{o}{.}\PY{n}{ones}\PY{p}{(}\PY{p}{(}\PY{l+m+mi}{2}\PY{p}{,} \PY{l+m+mi}{3}\PY{p}{)}\PY{p}{)}\PY{p}{,}
    \PY{l+s+s2}{\PYZdq{}}\PY{l+s+s2}{empty}\PY{l+s+s2}{\PYZdq{}}\PY{p}{:} \PY{n}{np}\PY{o}{.}\PY{n}{empty}\PY{p}{(}\PY{p}{(}\PY{l+m+mi}{2}\PY{p}{,} \PY{l+m+mi}{3}\PY{p}{)}\PY{p}{)}\PY{p}{,}
\PY{p}{\PYZcb{}}

\PY{n}{examples}
\end{Verbatim}
\end{tcolorbox}

            \begin{tcolorbox}[breakable, size=fbox, boxrule=.5pt, pad at break*=1mm, opacityfill=0]
\prompt{Out}{outcolor}{8}{\boxspacing}
\begin{Verbatim}[commandchars=\\\{\}]
\{'array': array([1, 2, 3]),
 'arange': array([0, 2, 4, 6, 8]),
 'linspace': array([0.  , 0.25, 0.5 , 0.75, 1.  ]),
 'zeros': array([[0., 0., 0.],
        [0., 0., 0.]]),
 'ones': array([[1., 1., 1.],
        [1., 1., 1.]]),
 'empty': array([[0., 0., 0.],
        [0., 0., 0.]])\}
\end{Verbatim}
\end{tcolorbox}
        
    \paragraph{ndarray
的内存占用}\label{ndarray-ux7684ux5185ux5b58ux5360ux7528}

我们用一个简单的例子介绍 ndarray 在存储数值上对比 \texttt{List} 的优势,
假设我们比较一百万个整数:

    \begin{tcolorbox}[breakable, size=fbox, boxrule=1pt, pad at break*=1mm,colback=cellbackground, colframe=cellborder]
\prompt{In}{incolor}{20}{\boxspacing}
\begin{Verbatim}[commandchars=\\\{\}]
\PY{n}{python\PYZus{}list} \PY{o}{=} \PY{n+nb}{list}\PY{p}{(}\PY{n+nb}{range}\PY{p}{(}\PY{l+m+mi}{1\PYZus{}000\PYZus{}000}\PY{p}{)}\PY{p}{)}
\PY{n}{numpy\PYZus{}array} \PY{o}{=} \PY{n}{np}\PY{o}{.}\PY{n}{arange}\PY{p}{(}\PY{l+m+mi}{1\PYZus{}000\PYZus{}000}\PY{p}{,} \PY{n}{dtype}\PY{o}{=}\PY{n}{np}\PY{o}{.}\PY{n}{int64}\PY{p}{)}

\PY{n}{python\PYZus{}approx\PYZus{}bytes} \PY{o}{=} \PY{p}{(}
    \PY{n}{sys}\PY{o}{.}\PY{n}{getsizeof}\PY{p}{(}\PY{n}{python\PYZus{}list}\PY{p}{)}
    \PY{o}{+} \PY{n+nb}{sum}\PY{p}{(}\PY{n}{sys}\PY{o}{.}\PY{n}{getsizeof}\PY{p}{(}\PY{n}{x}\PY{p}{)} \PY{k}{for} \PY{n}{x} \PY{o+ow}{in} \PY{n}{python\PYZus{}list}\PY{p}{)}
\PY{p}{)}

\PY{n+nb}{print}\PY{p}{(}
    \PY{l+s+s2}{\PYZdq{}}\PY{l+s+s2}{Python list 近似占用:}\PY{l+s+s2}{\PYZdq{}}\PY{p}{,}
    \PY{l+s+sa}{f}\PY{l+s+s2}{\PYZdq{}}\PY{l+s+si}{\PYZob{}}\PY{n}{python\PYZus{}approx\PYZus{}bytes}\PY{+w}{ }\PY{o}{/}\PY{+w}{ }\PY{l+m+mi}{1024}\PY{o}{*}\PY{o}{*}\PY{l+m+mi}{2}\PY{l+s+si}{:}\PY{l+s+s2}{,.2f}\PY{l+s+si}{\PYZcb{}}\PY{l+s+s2}{ MB}\PY{l+s+s2}{\PYZdq{}}
\PY{p}{)}

\PY{n+nb}{print}\PY{p}{(}
    \PY{l+s+s2}{\PYZdq{}}\PY{l+s+s2}{NumPy 元素缓冲区:}\PY{l+s+s2}{\PYZdq{}}\PY{p}{,}
    \PY{l+s+sa}{f}\PY{l+s+s2}{\PYZdq{}}\PY{l+s+si}{\PYZob{}}\PY{n}{numpy\PYZus{}array}\PY{o}{.}\PY{n}{nbytes}\PY{+w}{ }\PY{o}{/}\PY{+w}{ }\PY{l+m+mi}{1024}\PY{o}{*}\PY{o}{*}\PY{l+m+mi}{2}\PY{l+s+si}{:}\PY{l+s+s2}{,.2f}\PY{l+s+si}{\PYZcb{}}\PY{l+s+s2}{ MB}\PY{l+s+s2}{\PYZdq{}}
\PY{p}{)}
\end{Verbatim}
\end{tcolorbox}

    \begin{Verbatim}[commandchars=\\\{\}]
Python list 近似占用: 34.33 MB
NumPy 元素缓冲区: 7.63 MB
    \end{Verbatim}

    这里 Python 的统计是 CPython 对象大小近似, 而 \texttt{ndarray.nbytes}
统计的是元素缓冲区本身, 两者并不是完全相同口径; 这个实验的重点是理解:

\begin{quote}
\textbf{Python list 存的是一组 Python 对象引用, 而 NumPy
可以把同类型数值紧密存储.}
\end{quote}

但是不要为了``NumPy 看起来高级''把所有东西都放进 NumPy.

例如:

\begin{Shaded}
\begin{Highlighting}[]
\NormalTok{np.array(}
\NormalTok{    [}
\NormalTok{        \{}\StringTok{"name"}\NormalTok{: }\StringTok{"Alice"}\NormalTok{\},}
\NormalTok{        \{}\StringTok{"name"}\NormalTok{: }\StringTok{"Bob"}\NormalTok{\},}
\NormalTok{    ],}
\NormalTok{    dtype}\OperatorTok{=}\BuiltInTok{object}\NormalTok{,}
\NormalTok{)}
\end{Highlighting}
\end{Shaded}

此时底层仍然绕不开 Python object, 大量 NumPy 数值计算优势会消失.

    \paragraph{ndarray 的维度 Axis}\label{ndarray-ux7684ux7ef4ux5ea6-axis}

    在 ndarray 中,
\textbf{\texttt{axis}(轴)的本质是多维数组维度的``索引编号''与``遍历方向''}.
它定义了多维数据在逻辑结构中的组织层级,
以及算子在内存中沿哪个方向进行遍历或折叠.

\textbf{1. \texttt{axis} 的本质}

\begin{itemize}
\tightlist
\item
  \textbf{维度的编号:} 一个 \(N\) 维数组拥有从 \texttt{0} 到 \(N-1\)
  的轴(支持负数索引, \texttt{-1} 代表最后一维).
\item
  \textbf{与 \texttt{shape} 的直接映射:} \texttt{ndarray.shape}
  为一个元组 \((d_0, d_1, \dots, d_{N-1})\), 其中 \texttt{axis=i}
  对应的就是维度大小 \(d_i\).
\item
  \textbf{与嵌套列表层级的映射:} 最外层的中括号对应 \texttt{axis=0},
  每往里深入一层括号, \texttt{axis} 编号加 1.
\item
  \textbf{底层内存步长(Strides)的具象化:} 在默认的 C-order 连续内存中,
  \texttt{axis=0} 跨度最大(跳过一整个切片/行), \texttt{axis=-1}
  跨度最小(内存地址连续相邻).
\end{itemize}

\textbf{2. 指定 \texttt{axis} 到底意味着什么?}

理解 \texttt{axis} 计算最关键的一句话:\textbf{``沿着指定的 \texttt{axis}
方向移动并折叠, 消除该维度''}.

\textbf{3. \texttt{axis} 的核心作用}

\begin{itemize}
\tightlist
\item
  \textbf{统计与聚合(Reduction):} 控制
  \texttt{sum}、\texttt{mean}、\texttt{max}、\texttt{std}
  等函数的计算方向, 决定保留哪些维度、压缩哪些维度.
\item
  \textbf{拼接与扩展(Concatenation \& Stacking):}

  \begin{itemize}
  \tightlist
  \item
    \texttt{np.concatenate({[}a,\ b{]},\ axis=0)}:沿垂直方向拼接(增加行数).
  \item
    \texttt{np.concatenate({[}a,\ b{]},\ axis=1)}:沿水平方向拼接(增加列数).
  \end{itemize}
\item
  \textbf{排序与累加(Sorting \& Scanning):} 如
  \texttt{np.sort(arr,\ axis=1)} 仅对每一行内部的元素排序,
  行与行之间互不干扰.
\item
  \textbf{维度转换与重排(Transposition):} 在图像处理(如深度学习中
  \texttt{(H,\ W,\ C)} 转 \texttt{(C,\ H,\ W)})中, 通过
  \texttt{np.transpose(arr,\ axes=(2,\ 0,\ 1))} 重新排列各轴顺序.
\end{itemize}

    \paragraph{View、Copy 与 Strides}\label{viewcopy-ux4e0e-strides}

    在 ndarray 中,
\textbf{\texttt{Strides}(步长)是连接高维逻辑坐标与底层一维内存的``寻址导航器'',
而 \texttt{View}(视图)与
\texttt{Copy}(副本)是由此衍生的两种内存管理策略}. NumPy
之所以能在处理海量数据时保持极高的计算性能,
本质就在于将``数据实体''与``元数据(Metadata)''彻底解耦.

\textbf{为什么必须放在一起讲?}

\texttt{View} 的物理实现完全依赖于 \texttt{Strides} 的重新映射; 掌握了
\texttt{Strides} 的线性寻址逻辑,
就能一眼看穿为什么基础切片可以实现零拷贝(\(O(1)\) 复杂度),
而花式索引却必须退化为内存深拷贝(\(O(N)\) 复杂度).

\textbf{1. 本质}

\begin{itemize}
\tightlist
\item
  \textbf{双层架构解耦:} \texttt{ndarray} 由两部分构成:
  保存纯字节流的\textbf{连续数据缓冲区(Data
  Buffer)*\emph{与保存描述信息的*}元数据头(Metadata Header)}(包含
  \texttt{shape}、\texttt{dtype}、\texttt{strides}、\texttt{data\ pointer}、\texttt{base}
  等).
\item
  \textbf{\texttt{Strides}(步长)的本质:} 一个元组
  \((s_0, s_1, \dots, s_{N-1})\), 表示在内存字节流中沿 \texttt{axis=i}
  移动一个索引单位需要跨越的\textbf{物理字节数(Bytes)}.
\item
  \textbf{\texttt{View}(视图)的本质:} 创建了一个全新的元数据头, 其
  \texttt{data\ pointer} 仍指向\textbf{同一个底层内存缓冲区}.
  创建时间复杂度为 \(O(1)\),
  对视图的数据修改会直接同步到原数组(产生副作用).
\item
  \textbf{\texttt{Copy}(副本)的本质:}
  在堆内存中开辟了一块全新的独立数据缓冲区, 并将数据完整复制过去.
  时间复杂度为 \(O(N)\), 与原数组在物理内存上完全隔离.
\end{itemize}

\textbf{2. 原理}

\begin{itemize}
\item
  \textbf{内存线性寻址公式:} 任意多维逻辑索引
  \((i_0, i_1, \dots, i_{N-1})\) 对应的底层字节偏移量为:

  \[\text{Byte Offset} = \sum_{k=0}^{N-1} (i_k \times s_k)\]
\item
  \textbf{View vs Copy 的判定准则:}

  \begin{itemize}
  \tightlist
  \item
    \textbf{产生 View:}
    只要操作后的新数据子集能够通过\textbf{单组等差步长(Strides)与基地址偏移}完全表达(如切片
    \texttt{arr{[}::2{]}}、转置 \texttt{arr.T}、维度压缩
    \texttt{arr.squeeze()}), NumPy 就仅修改元数据返回 View.
  \item
    \textbf{强制 Copy:}
    一旦目标元素的内存分布\textbf{无法用固定的线性步长表达}(如花式索引
    \texttt{arr{[}{[}0,\ 2{]}{]}}、布尔掩码
    \texttt{arr{[}arr\ \textgreater{}\ 5{]}}), 或显式调用
    \texttt{.copy()}, NumPy 必须分配新内存并产生 Copy.
  \end{itemize}
\end{itemize}

\textbf{3. 核心作用}

\begin{itemize}
\tightlist
\item
  \textbf{极速性能(Zero-Copy):}
  切片、翻转、转置等操作耗时与数组大小完全无关(均在纳秒级完成),
  避免大规模内存申请与数据搬运.
\item
  \textbf{内存安全控制:} 显式调用 \texttt{.copy()}
  切断与原数组的指针绑定,
  防止多线程计算或数据预处理流水线中的意外内存污染.
\item
  \textbf{内存连续性(Contiguity)管理:} ndarray 的数据是否按照 C
  语言/行优先(row-major)的方式连续存储, 确保调用底层 BLAS/LAPACK 或
  C/C++ 扩展时的 SIMD 向量化加速效率.
\item
  \textbf{高级数据重组(Strided Tricks):} 通过手动操控 Strides,
  无需分配额外内存即可实现图像滑动窗口(Sliding
  Window)、重叠切片或卷积核展开.
\end{itemize}

    \subsubsection{ndarray 的用法}\label{ndarray-ux7684ux7528ux6cd5}

    \paragraph{广播机制(Broadcasting)}\label{ux5e7fux64adux673aux5236broadcasting}

    在 ndarray 的算术运算中, \textbf{\texttt{Broadcasting}(广播机制)是 NumPy
在不复制数据的前提下, 自动拉伸和对齐不同形状数组的``隐式虚拟扩展引擎''}.
它是向量化计算(Vectorization)的核心支柱, 让具有不同形状的张量能够在底层
C 语言级别无缝进行逐元素(Element-wise)运算.

\textbf{1. 本质}

\begin{itemize}
\tightlist
\item
  \textbf{零内存消耗的虚拟扩展:}
  广播并不是在物理内存中真实复制数据填充数组,
  而是通过\textbf{将对应拉伸维度的步长(Strides)设为 \texttt{0}}. 步长为
  0 意味着逻辑索引在递增, 但底层指针停留在同一物理内存地址上, 实现
  \(O(1)\) 内存开销的``逻辑克隆''.
\item
  \textbf{双重规约机制:} 广播定义了严格的维度靠右对齐与长度兼容法则,
  将多维数组间的算术运算、比较运算和赋值操作彻底标准化.
\end{itemize}

\textbf{2. 原理(对齐双法则)}

广播在底层严格按照以下两个阶段执行判断与形状推导:

\begin{itemize}
\tightlist
\item
  \textbf{法则一: 维度补齐(靠右对齐, 左侧补 1)}

  \begin{itemize}
  \tightlist
  \item
    如果两个数组的维度数(\texttt{ndim})不同, NumPy 会在维度较少的数组
    \texttt{shape} \textbf{左侧补 1}, 直到两者维度数量一致.
  \item
    \emph{例:} \texttt{shape=(3,)} 与 \texttt{shape=(4,\ 3)} 对齐时,
    前者先被补齐为 \texttt{(1,\ 3)}.
  \end{itemize}
\item
  \textbf{法则二: 长度兼容(等于自身或等于 1)}

  \begin{itemize}
  \tightlist
  \item
    从最后一个维度(Trailing Dimension, 即最右侧轴)开始向前逐轴比对.
    对于每个轴, 两者的长度必须满足以下条件之一, 否则抛出
    \texttt{ValueError}:

    \begin{enumerate}
    \def\labelenumi{\arabic{enumi}.}
    \tightlist
    \item
      两个维度的长度相等;
    \item
      其中一个维度的长度为
      \texttt{1}(该维度会被虚拟拉伸到与另一个维度相同).
    \end{enumerate}
  \end{itemize}
\item
  \textbf{输出形状判定:} 最终输出结果在每个轴上的长度,
  等于参与运算各数组在该轴长度的\textbf{最大值}:
  \(\max(d_{1,i}, d_{2,i})\).
\end{itemize}

\textbf{3. 核心作用}

\begin{itemize}
\tightlist
\item
  \textbf{消除显式 Python 循环:} 避免书写多重嵌套循环, 将计算下沉至底层
  C 连续遍历, 极大激发 CPU 的 SIMD 向量化指令集性能.
\item
  \textbf{内存极致优化:} 彻底淘汰 \texttt{np.tile} 或 \texttt{np.repeat}
  等物理深拷贝复制数据的低效做法.
\item
  \textbf{高维特征与批处理运算:} 在机器学习与图像处理中,
  轻松实现全局特征去中心化、样本按权重缩放、网格生成(Meshgrid)以及成对距离矩阵(Pairwise
  Distance)计算.
\end{itemize}

    \paragraph{Vectorization 与
UFunc(通用函数)}\label{vectorization-ux4e0e-ufuncux901aux7528ux51fdux6570}

    Vectorization(向量化)是 NumPy 高效计算的编程思想/使用方式, 而
UFunc(通用函数)是 ndarray 实现高效逐元素计算的核心机制之一.
它们结合``连续内存 + 固定 dtype + C 层循环 + CPU
向量化/SIMD''等一整套机制共同构建了 Numpy 的高性能计算体系.
\textbf{\texttt{Vectorization}(向量化)是将原本由 Python
解释器执行的逐元素标量循环, 下沉至底层 C 语言连续内存块与 CPU SIMD
硬件指令的``去循环执行模式''; 而 \texttt{UFunc}(Universal Function,
通用函数)则是封装了这种高效 C
语言内核、并自带广播与高阶聚合方法的``可调用对象引擎''}.

\textbf{1. 本质}

\begin{itemize}
\tightlist
\item
  \textbf{消除解释器开销(Bypassing Interpreter Overhead):} 摆脱 Python
  在每次循环中反复进行的动态类型检查(Dynamic Type
  Resolution)、对象封包解包(Boxing/Unboxing)以及引用计数维护,
  将计算交由预编译的 C 循环指针递增执行.
\item
  \textbf{UFunc 对象的物理本质:} 一个内建了多重类型分发表(Type Dispatch
  Table)与连续内存迭代器(\texttt{NpyIter})的 C
  结构体对象(\texttt{PyUFuncObject}), 它同时管理着一元算子(如
  \texttt{sin}、\texttt{exp})与二元算子(如
  \texttt{add}、\texttt{multiply}).
\item
  \textbf{硬件级并行(SIMD 并发):} 向量化使连续内存数据能够直接加载进现代
  CPU 的向量寄存器(如 AVX-512、AVX2、ARM NEON),
  实现单指令周期内对多个浮点数的并发吞吐(SIMD, Single Instruction
  Multiple Data).
\end{itemize}

\textbf{2. 原理}

\begin{itemize}
\tightlist
\item
  \textbf{类型分发(Type Resolution):} 传入数组时, UFunc 根据输入的
  \texttt{dtype} 匹配最精确的底层 C 函数指针(如 \texttt{int64}
  走整数加法指令, \texttt{float64} 走双精度浮点流水线).
\item
  \textbf{天然嵌入广播:} 任何多元 UFunc 都会自动调用广播协议,
  无需额外处理维度即可自动扩展输入形状.
\item
  \textbf{方法派生矩阵(UFunc Methods):} 任何\textbf{二元通用函数}(Binary
  UFunc, 如 \texttt{np.add}、\texttt{np.multiply}、\texttt{np.maximum}
  等)都自动派生出 5 个高阶方法,
  直接将简单的逐元素算子转化为强大的矩阵与序列计算引擎:
\end{itemize}

\begin{enumerate}
\def\labelenumi{\arabic{enumi}.}
\tightlist
\item
  \texttt{.reduce()}:沿指定轴折叠压扁(累加/累乘/最值筛选);
\item
  \texttt{.accumulate()}:沿指定轴保存中间累积状态(前缀和/前缀积);
\item
  \texttt{.outer()}:两输入做笛卡尔积全组合外积运算;
\item
  \texttt{.at()}:无缓冲的就地离散索引更新(支持重复索引累加);
\item
  \texttt{.reduceat()}:按指定的索引切片边界分段折叠汇总.
\end{enumerate}

\textbf{3. 核心作用}

\begin{itemize}
\tightlist
\item
  \textbf{极速性能吞吐:} 通常比原生 Python 的 \texttt{for}
  循环与列表推导式带来 50 至 500 倍的性能提升.
\item
  \textbf{高阶函数式表达:} 用 \texttt{.outer()} 或 \texttt{.reduce()}
  替代复杂的多层循环与动态列表收集, 大幅精简工程代码.
\item
  \textbf{解决原地累加数据冲突(Race Hazard):} \texttt{ufunc.at}
  解决了常规花式索引(Fancy
  Indexing)在遇到重复索引时因写入缓冲(Buffering)导致累加丢失的致命问题.
\item
  \textbf{灵活的接口封装:} 借助 \texttt{np.vectorize} 与
  \texttt{np.frompyfunc}, 可将任意纯 Python
  业务函数快速封装为兼容广播特性的伪向量化对象.
\end{itemize}

    \subsection{Pandas:}\label{pandas}

    我通过前面抛出的几个问题来解释 Pandas 的实际用法.

{\def\LTcaptype{none} % do not increment counter
\begin{longtable}[]{@{}
  >{\centering\arraybackslash}p{(\linewidth - 4\tabcolsep) * \real{0.3333}}
  >{\centering\arraybackslash}p{(\linewidth - 4\tabcolsep) * \real{0.3333}}
  >{\centering\arraybackslash}p{(\linewidth - 4\tabcolsep) * \real{0.3333}}@{}}
\toprule\noalign{}
\begin{minipage}[b]{\linewidth}\centering
问题
\end{minipage} & \begin{minipage}[b]{\linewidth}\centering
业务问题
\end{minipage} & \begin{minipage}[b]{\linewidth}\centering
逐步引出的技术
\end{minipage} \\
\midrule\noalign{}
\endhead
\bottomrule\noalign{}
\endlastfoot
\textbf{Q1} & 这 200 万条数据可靠吗? &
\texttt{dtype}、缺失值、mask、\texttt{loc}、\texttt{query} \\
\textbf{Q2} & 一天中什么时候最忙? &
\texttt{datetime}、\texttt{groupby}、\texttt{agg} \\
\textbf{Q3} & 什么时间段的流水和单位运营时间收入最高? &
向量化、\texttt{assign}、\texttt{groupby}、派生指标 \\
\textbf{Q4} & 哪些上下车区域最热门? &
\texttt{value\_counts}、\texttt{merge}、连接校验 \\
\textbf{Q5} & 哪些订单和时间窗口可能异常? &
\texttt{where}、\texttt{select}、\texttt{quantile}、\texttt{transform}、\texttt{resample}、\texttt{rolling} \\
\textbf{Q6} & 为什么同一个分析可以相差一个数量级甚至更多? &
\texttt{Python\ loop}、\texttt{apply}、\texttt{NumPy}、\texttt{Numba}、\texttt{eval/query}、内存优化 \\
\textbf{Q7} & 如何把分析脚本变成公司可以重复使用的工具? &
\texttt{pipe}、纯函数、\texttt{Accessor}、领域 API \\
\end{longtable}
}

基于问题 Q1-Q6 介绍 Pandas 的核心设计思想与用法.

    \subsubsection{Pandas
的设计哲学}\label{pandas-ux7684ux8bbeux8ba1ux54f2ux5b66}

    \paragraph{Series、DataFrame 和
Index}\label{seriesdataframe-ux548c-index}

    从 NumPy 到 Pandas 的设计思想, 其实是在回答一个非常现实的问题:

\begin{quote}
\textbf{当数据不再只是``一坨数字'',
而是带着字段名、业务主键、缺失值、时间戳和连接关系的表格时,
我们还能够继续只靠 \texttt{ndarray} 吗?}
\end{quote}

答案通常是不够.

NumPy 擅长的是 \textbf{同构数值数组} 的高性能计算; 而 Pandas 擅长的是
\textbf{带业务语义的异构表格数据}
的组织、筛选、对齐、聚合、连接与时间序列分析.

也就是说:

\begin{Shaded}
\begin{Highlighting}[]
\NormalTok{NumPy 关心的是:}
\NormalTok{这块连续内存怎样更快地算? }

\NormalTok{Pandas 关心的是:}
\NormalTok{这张表里的哪一列代表时间? }
\NormalTok{哪一列代表金额? }
\NormalTok{哪些行缺失? }
\NormalTok{两张表如何按主键对齐? }
\NormalTok{分组后如何再把结果贴回原表? }
\end{Highlighting}
\end{Shaded}

这就是 Pandas 的设计出发点:

\begin{quote}
\textbf{在 NumPy 的高性能数组之上,
增加标签、索引、缺失值语义、关系代数和时间序列能力,
让``数据分析代码''既能跑得够快, 又能写得像业务逻辑.}
\end{quote}

    \paragraph{Pandas 的三块根基:
Series、DataFrame、Index}\label{pandas-ux7684ux4e09ux5757ux6839ux57fa-seriesdataframeindex}

    \subparagraph{\texorpdfstring{\texttt{Series}
的本质}{Series 的本质}}\label{series-ux7684ux672cux8d28}

    \texttt{Series} 可以理解为:

\begin{quote}
\textbf{一列带标签的一维数组 = values + index}
\end{quote}

它不是简单的一维 list, 也不是裸的 ndarray,
而是``数值/对象数组''与``行标签''绑定后的结果.

\begin{Shaded}
\begin{Highlighting}[]
\NormalTok{s }\OperatorTok{=}\NormalTok{ pd.Series([}\FloatTok{10.0}\NormalTok{, }\FloatTok{20.0}\NormalTok{, }\FloatTok{30.0}\NormalTok{], index}\OperatorTok{=}\NormalTok{[}\StringTok{"A"}\NormalTok{, }\StringTok{"B"}\NormalTok{, }\StringTok{"C"}\NormalTok{], name}\OperatorTok{=}\StringTok{"fare"}\NormalTok{)}

\BuiltInTok{print}\NormalTok{(}\StringTok{"Series 内容:}\CharTok{\textbackslash{}n}\StringTok{"}\NormalTok{, s)}
\BuiltInTok{print}\NormalTok{(}\StringTok{"values:"}\NormalTok{, s.values)}
\BuiltInTok{print}\NormalTok{(}\StringTok{"index :"}\NormalTok{, s.index)}
\BuiltInTok{print}\NormalTok{(}\StringTok{"name  :"}\NormalTok{, s.name)}
\BuiltInTok{print}\NormalTok{(}\StringTok{"dtype :"}\NormalTok{, s.dtype)}
\end{Highlighting}
\end{Shaded}

这意味着 \texttt{Series} 的运算不是单纯按位置发生,
而是可以按标签自动对齐:

\begin{Shaded}
\begin{Highlighting}[]
\NormalTok{s1 }\OperatorTok{=}\NormalTok{ pd.Series([}\DecValTok{10}\NormalTok{, }\DecValTok{20}\NormalTok{, }\DecValTok{30}\NormalTok{], index}\OperatorTok{=}\NormalTok{[}\StringTok{"A"}\NormalTok{, }\StringTok{"B"}\NormalTok{, }\StringTok{"C"}\NormalTok{])}
\NormalTok{s2 }\OperatorTok{=}\NormalTok{ pd.Series([}\DecValTok{1}\NormalTok{, }\DecValTok{2}\NormalTok{, }\DecValTok{3}\NormalTok{], index}\OperatorTok{=}\NormalTok{[}\StringTok{"B"}\NormalTok{, }\StringTok{"C"}\NormalTok{, }\StringTok{"D"}\NormalTok{])}

\BuiltInTok{print}\NormalTok{(}\StringTok{"按索引对齐相加的结果:}\CharTok{\textbackslash{}n}\StringTok{"}\NormalTok{, s1 }\OperatorTok{+}\NormalTok{ s2)}
\end{Highlighting}
\end{Shaded}

输出会是:

\begin{Shaded}
\begin{Highlighting}[]
\NormalTok{A     NaN}
\NormalTok{B    21.0}
\NormalTok{C    32.0}
\NormalTok{D     NaN}
\NormalTok{dtype: float64}
\end{Highlighting}
\end{Shaded}

因为 Pandas 默认认为:

\begin{quote}
\textbf{业务上``B''只能和``B''对齐, ``C''只能和``C''对齐,
没匹配到的就应当显式表现为缺失值 \texttt{NaN}.}
\end{quote}

这就是 Pandas 跟 NumPy 在抽象层次上最本质的区别之一.

    \subparagraph{\texorpdfstring{\texttt{DataFrame}
的本质}{DataFrame 的本质}}\label{dataframe-ux7684ux672cux8d28}

    \texttt{DataFrame} 可以理解为:

\begin{quote}
\textbf{多个共享同一套行索引的 \texttt{Series} 组成的二维表结构}
\end{quote}

它既可以看作``带列名和行索引的二维表'',
也可以看作``按列组织的一组一维数组''.

\begin{Shaded}
\begin{Highlighting}[]
\NormalTok{df\_demo }\OperatorTok{=}\NormalTok{ pd.DataFrame(}
\NormalTok{    \{}
        \StringTok{"fare\_amount"}\NormalTok{: [}\FloatTok{8.5}\NormalTok{, }\FloatTok{12.0}\NormalTok{, }\FloatTok{6.0}\NormalTok{],}
        \StringTok{"passenger\_count"}\NormalTok{: [}\DecValTok{1}\NormalTok{, }\DecValTok{2}\NormalTok{, }\DecValTok{1}\NormalTok{],}
        \StringTok{"pickup\_hour"}\NormalTok{: [}\DecValTok{8}\NormalTok{, }\DecValTok{9}\NormalTok{, }\DecValTok{8}\NormalTok{],}
\NormalTok{    \},}
\NormalTok{    index}\OperatorTok{=}\NormalTok{[}\StringTok{"trip\_1"}\NormalTok{, }\StringTok{"trip\_2"}\NormalTok{, }\StringTok{"trip\_3"}\NormalTok{],}
\NormalTok{)}

\BuiltInTok{print}\NormalTok{(}\StringTok{"DataFrame 内容:}\CharTok{\textbackslash{}n}\StringTok{"}\NormalTok{, df\_demo)}
\BuiltInTok{print}\NormalTok{(}\StringTok{"index   :"}\NormalTok{, df\_demo.index)}
\BuiltInTok{print}\NormalTok{(}\StringTok{"columns :"}\NormalTok{, df\_demo.columns)}
\BuiltInTok{print}\NormalTok{(}\StringTok{"各列 dtype:"}\NormalTok{)}
\BuiltInTok{print}\NormalTok{(df\_demo.dtypes)}
\end{Highlighting}
\end{Shaded}

从物理角度看, \texttt{DataFrame} 更接近 \textbf{列式存储}
而不是``一个个行对象''.

这件事非常关键, 因为数据分析的大量操作都是:

\begin{itemize}
\tightlist
\item
  对某一列做过滤;
\item
  对某几列做数值计算;
\item
  对某一列做分组统计;
\item
  对时间列做重采样;
\item
  对主键列做连接.
\end{itemize}

列式组织意味着 Pandas 可以更高效地把这些操作下压到 NumPy
或扩展数组层面执行.

    \subparagraph{\texorpdfstring{\texttt{Index}
\textbar 的本质}{Index \textbar 的本质}}\label{index-ux7684ux672cux8d28}

    很多初学者把 \texttt{Index} 理解成``行号'', 这是不够的.

\texttt{Index} 的真正角色是:

\begin{quote}
\textbf{Pandas 中``轴标签''的统一抽象,
用于定位、对齐、分组、切片、连接和时间序列操作.}
\end{quote}

它既可以是默认的 \texttt{RangeIndex}, 也可以是:

\begin{itemize}
\tightlist
\item
  字符串主键组成的 \texttt{Index}
\item
  时间戳组成的 \texttt{DatetimeIndex}
\item
  多层标签组成的 \texttt{MultiIndex}
\end{itemize}

例如:

\begin{Shaded}
\begin{Highlighting}[]
\NormalTok{df\_idx }\OperatorTok{=}\NormalTok{ pd.DataFrame(}
\NormalTok{    \{}
        \StringTok{"fare\_amount"}\NormalTok{: [}\DecValTok{10}\NormalTok{, }\DecValTok{12}\NormalTok{, }\DecValTok{8}\NormalTok{],}
        \StringTok{"passenger\_count"}\NormalTok{: [}\DecValTok{1}\NormalTok{, }\DecValTok{2}\NormalTok{, }\DecValTok{1}\NormalTok{],}
\NormalTok{    \},}
\NormalTok{    index}\OperatorTok{=}\NormalTok{[}\StringTok{"order\_001"}\NormalTok{, }\StringTok{"order\_002"}\NormalTok{, }\StringTok{"order\_003"}\NormalTok{],}
\NormalTok{)}

\BuiltInTok{print}\NormalTok{(}\StringTok{"按标签定位 order\_002 这一行:}\CharTok{\textbackslash{}n}\StringTok{"}\NormalTok{, df\_idx.loc[}\StringTok{"order\_002"}\NormalTok{])}
\end{Highlighting}
\end{Shaded}

如果把时间列设置为索引:

\begin{Shaded}
\begin{Highlighting}[]
\NormalTok{df\_time }\OperatorTok{=}\NormalTok{ pd.DataFrame(}
\NormalTok{    \{}
        \StringTok{"pickup\_count"}\NormalTok{: [}\DecValTok{120}\NormalTok{, }\DecValTok{150}\NormalTok{, }\DecValTok{180}\NormalTok{],}
\NormalTok{    \},}
\NormalTok{    index}\OperatorTok{=}\NormalTok{pd.to\_datetime([}
        \StringTok{"2024{-}01{-}01 08:00:00"}\NormalTok{,}
        \StringTok{"2024{-}01{-}01 09:00:00"}\NormalTok{,}
        \StringTok{"2024{-}01{-}01 10:00:00"}\NormalTok{,}
\NormalTok{    ]),}
\NormalTok{)}

\BuiltInTok{print}\NormalTok{(}\StringTok{"按时间区间切片(08:00 至 09:00):}\CharTok{\textbackslash{}n}\StringTok{"}\NormalTok{, df\_time.loc[}\StringTok{"2024{-}01{-}01 08:00:00"}\NormalTok{:}\StringTok{"2024{-}01{-}01 09:00:00"}\NormalTok{])}
\end{Highlighting}
\end{Shaded}

此时 Pandas 就能直接支持基于时间的切片、重采样与滑窗.

    \paragraph{Pandas
的设计哲学总结}\label{pandas-ux7684ux8bbeux8ba1ux54f2ux5b66ux603bux7ed3}

    把上面三者放在一起看, Pandas 的核心设计哲学可以概括为四句话:

\textbf{1. 在数组之上增加标签语义}

NumPy 的世界里, 重点是位置; Pandas 的世界里, 重点是字段名和标签.

\textbf{2. 在高性能基础上容纳异构数据}

现实业务表不是纯 \texttt{float64} 矩阵,
而是金额、字符串、时间、布尔、缺失值混合在一起.

\textbf{3. 把关系代数和时间序列作为一等公民}

\texttt{groupby}、\texttt{merge}、\texttt{resample}、\texttt{rolling} 在
Pandas 里不是边角料, 而是最核心的分析算子.

\textbf{4. 让分析代码表达业务意图而不是遍历细节}

不是反复写:

\begin{Shaded}
\begin{Highlighting}[]
\ControlFlowTok{for}\NormalTok{ row }\KeywordTok{in}\NormalTok{ rows:}
\NormalTok{    ...}
\end{Highlighting}
\end{Shaded}

而是写成:

\begin{Shaded}
\begin{Highlighting}[]
\NormalTok{(}
\NormalTok{    df}
\NormalTok{    .query(}\StringTok{"total\_amount \textgreater{} 0"}\NormalTok{)}
\NormalTok{    .assign(pickup\_hour}\OperatorTok{=}\KeywordTok{lambda}\NormalTok{ x: x[}\StringTok{"tpep\_pickup\_datetime"}\NormalTok{].dt.hour)}
\NormalTok{    .groupby(}\StringTok{"pickup\_hour"}\NormalTok{)}
\NormalTok{    .agg(order\_count}\OperatorTok{=}\NormalTok{(}\StringTok{"total\_amount"}\NormalTok{, }\StringTok{"size"}\NormalTok{))}
\NormalTok{)}
\end{Highlighting}
\end{Shaded}

这类代码表达的是``业务分析动作'', 而不是``解释器怎样逐行跑''.

    \subsubsection{Q1 :这 200
万条数据可靠吗?}\label{q1-ux8fd9-200-ux4e07ux6761ux6570ux636eux53efux9760ux5417}

    这个问题看起来很朴素, 但实际上是所有分析的起点.

如果数据本身就不可信:

\begin{itemize}
\tightlist
\item
  类型错了;
\item
  时间还没转成真正的时间戳;
\item
  缺失值没有识别;
\item
  坐标出现 0 或离谱异常值;
\item
  金额出现负数或极端脏值;
\end{itemize}

那么后面所有均值、分组、建模、可视化都会建立在错误基础上.

所以 Q1 的真正目标不是``会几个 API'', 而是建立一个习惯:

\begin{quote}
\textbf{分析前先做数据可信度诊断(data validation / data quality check).}
\end{quote}

    \paragraph{这一单元要解决的
API}\label{ux8fd9ux4e00ux5355ux5143ux8981ux89e3ux51b3ux7684-api}

    {\def\LTcaptype{none} % do not increment counter
\begin{longtable}[]{@{}ll@{}}
\toprule\noalign{}
技术 & 在这里解决什么问题 \\
\midrule\noalign{}
\endhead
\bottomrule\noalign{}
\endlastfoot
\texttt{dtype} & 这一列到底是数值、字符串还是时间? \\
缺失值 & 哪些列有 \texttt{NaN} / \texttt{None} / 空洞? \\
\texttt{mask} & 如何批量构造``可信/不可信''的布尔条件? \\
\texttt{loc} & 如何按布尔条件精确筛出异常行? \\
\texttt{query} & 如何用接近 SQL 的写法进行业务过滤? \\
\end{longtable}
}

    \paragraph{本质}\label{ux672cux8d28}

    Q1 本质上做的是三件事:

\begin{enumerate}
\def\labelenumi{\arabic{enumi}.}
\tightlist
\item
  识别列的数据类型是否符合业务预期;
\item
  识别缺失值与异常值;
\item
  根据布尔规则把``可信数据子集''筛选出来.
\end{enumerate}

Pandas 在这里的价值是:

\begin{quote}
\textbf{把逐行检查, 变成整列向量化检查; 把``人工目测'',
变成``可重复执行的规则''.}
\end{quote}

    \paragraph{底层原理}\label{ux5e95ux5c42ux539fux7406}

    这一类操作底层主要依赖三件事:

\textbf{第一, 每一列都有明确的 \texttt{dtype} 与缺失值表示语义}

Pandas 会根据列类型选择不同的存储与计算路径. 数值列可以直接批量比较;
对象列通常更贵; 时间列转换后才能启用 \texttt{.dt}、时间切片和重采样能力.

\textbf{第二, 布尔掩码本质上是一列 \texttt{True/False}}

例如:

\begin{Shaded}
\begin{Highlighting}[]
\NormalTok{mask }\OperatorTok{=}\NormalTok{ df[}\StringTok{"fare\_amount"}\NormalTok{] }\OperatorTok{\textgreater{}} \DecValTok{0}
\end{Highlighting}
\end{Shaded}

这个 \texttt{mask} 的本质是一个与 DataFrame 行数等长的布尔
\texttt{Series}, 它像 NumPy 里的布尔索引一样, 但带有索引语义.

\textbf{第三, \texttt{loc} 与 \texttt{query}
都是在做``基于条件的行过滤''}

\begin{itemize}
\tightlist
\item
  \texttt{loc} 更显式, 更适合组合多个条件和精确选列;
\item
  \texttt{query} 更接近声明式 DSL, 可读性更强, 在复杂业务表达里很方便.
\end{itemize}

    \paragraph{数据分析中的作用与用法}\label{ux6570ux636eux5206ux6790ux4e2dux7684ux4f5cux7528ux4e0eux7528ux6cd5}

    在真实项目里, Q1 一般会固定形成一个``数据准入规则''阶段:

\begin{itemize}
\tightlist
\item
  先看 \texttt{df.dtypes}
\item
  再看 \texttt{df.isna().sum()}
\item
  再看描述统计 \texttt{describe()} 或分位数
\item
  再写出异常规则 \texttt{mask}
\item
  最后筛出 \texttt{clean\_df}
\end{itemize}

这一步非常像后端系统里的参数校验, 只不过对象从``一次请求''变成了``200
万条记录''.

如果把这道题真正落到 Pandas API 上, 可以按下面顺序思考:

\begin{enumerate}
\def\labelenumi{\arabic{enumi}.}
\tightlist
\item
  先用 \texttt{pd.read\_parquet}、\texttt{shape}、\texttt{dtypes}
  认清数据结构和字段类型;
\item
  再用 \texttt{isna}、数值比较、时间差计算等方式识别缺失值和异常值;
\item
  用多个布尔 \texttt{mask} 组合成``可信数据''的准入规则;
\item
  用 \texttt{loc} 抽查异常样本, 用 \texttt{query} 产出清洗后的可信子集;
\item
  最后用 \texttt{np.select} 把异常进一步分类成可统计、可追踪的标签.
\end{enumerate}

    \paragraph{真正需要记住什么?}\label{ux771fux6b63ux9700ux8981ux8bb0ux4f4fux4ec0ux4e48}

    不是``我会 \texttt{query} 和 \texttt{loc}''. 而是:

\begin{quote}
\textbf{Pandas 的第一价值,
是让数据质量检查变成整列、批量、可复现的规则系统, 而不是人工抽样拍脑袋.}
\end{quote}

    \subsubsection{Pandas: GroupBy、Transform、Merge 对应着问题
Q2～Q4}\label{pandas-groupbytransformmerge-ux5bf9ux5e94ux7740ux95eeux9898-q2q4}

    Q2: 一天中什么时候最忙?\\
Q3: 什么时间段的流水和单位运营时间收入最高?\\
Q4: 哪些上下车区域最热门?

Q2 到 Q4 其实对应了三类极其核心的分析动作:

\begin{itemize}
\tightlist
\item
  Q2:聚合统计
\item
  Q3:分组后再回填到原表
\item
  Q4:把不同来源的数据拼接起来
\end{itemize}

如果说 Q1 是``先把数据洗干净'', 那么 Q2～Q4 就是``开始真正做业务分析''.

    \paragraph{GroupBy、Transform、Merge}\label{groupbytransformmerge}

    \subparagraph{\texorpdfstring{\texttt{groupby}
的本质}{groupby 的本质}}\label{groupby-ux7684ux672cux8d28}

    这里的 \texttt{groupby} 和 MySQL 中的 \texttt{groupby} 十分类似,
它的经典模型是:

\begin{quote}
\textbf{Split - Apply - Combine}
\end{quote}

也就是:

\begin{enumerate}
\def\labelenumi{\arabic{enumi}.}
\tightlist
\item
  按某个键把数据拆成多个组;
\item
  对每个组做计算;
\item
  再把结果组合起来.
\end{enumerate}

例如``按小时统计订单量'':

\begin{Shaded}
\begin{Highlighting}[]
\NormalTok{df.groupby(}\StringTok{"pickup\_hour"}\NormalTok{).size()}
\end{Highlighting}
\end{Shaded}

    \subparagraph{\texorpdfstring{\texttt{transform}
的本质}{transform 的本质}}\label{transform-ux7684ux672cux8d28}

    \texttt{transform} 很容易和 \texttt{agg} 混淆.

二者的关键区别是:

\begin{itemize}
\tightlist
\item
  \texttt{agg} 会把每组压缩成一个或少数几个结果;
\item
  \texttt{transform} 会把每组算出的结果, \textbf{按原行数广播回去}.
\end{itemize}

也就是说:

\begin{quote}
\textbf{\texttt{transform} 适合做``分组特征回填'', 而 \texttt{agg}
适合做``分组汇总报表''.}
\end{quote}

    \subparagraph{\texorpdfstring{\texttt{merge}
的本质}{merge 的本质}}\label{merge-ux7684ux672cux8d28}

    \texttt{merge} 的本质就是表与表之间基于键的连接.

它在 Pandas 中扮演的角色, 基本等价于 SQL 里的 Join:

\begin{itemize}
\tightlist
\item
  \texttt{left}
\item
  \texttt{right}
\item
  \texttt{inner}
\item
  \texttt{outer}
\end{itemize}

而且 Pandas 额外强调:

\begin{quote}
\textbf{连接不仅是``拼起来'', 更重要的是``验证有没有拼错''.}
\end{quote}

这也是为什么 \texttt{validate=} 参数在教学里非常值得讲.

    \paragraph{核心特点: Index
对齐机制}\label{ux6838ux5fc3ux7279ux70b9-index-ux5bf9ux9f50ux673aux5236}

    Pandas 的底层灵魂是 \textbf{Index 对齐}, 这三个操作的核心差异,
本质上就是对 Index 的不同处理策略:

\begin{enumerate}
\def\labelenumi{\arabic{enumi}.}
\tightlist
\item
  \textbf{\texttt{groupby\ (+\ agg)}: 索引升格与压缩(降维)}
\end{enumerate}

\begin{itemize}
\tightlist
\item
  \textbf{特点}: 把指定的分组键提拔为\textbf{新的 Index}.
\item
  \textbf{机制}: 原表的 \(N\) 行被压缩为分组数 \(K\) 行,
  聚合结果与\textbf{新生成的 Group Index} 严格对齐.
\end{itemize}

\begin{enumerate}
\def\labelenumi{\arabic{enumi}.}
\setcounter{enumi}{1}
\tightlist
\item
  \textbf{\texttt{transform}: 严格保留原索引(保维)}
\end{enumerate}

\begin{itemize}
\tightlist
\item
  \textbf{特点}: 计算结果\textbf{强制与原表的原始 Index
  保持对齐}并原路广播.
\item
  \textbf{机制}: 输出始终保持 \(N\) 行, 无需遍历也无需手动关联,
  天然保证行索引与原表一行一对应.
\end{itemize}

\begin{enumerate}
\def\labelenumi{\arabic{enumi}.}
\setcounter{enumi}{2}
\tightlist
\item
  \textbf{\texttt{merge}: 跨表索引重构与映射(改维)}
\end{enumerate}

\begin{itemize}
\tightlist
\item
  \textbf{特点}: 基于键或 Index 寻找交集/并集,
  并\textbf{重构生成一套全新的 Index}.
\item
  \textbf{机制}: 处理跨表对齐, 根据匹配关系(1:1、1:N、N:N)对行和 Index
  进行保留、丢弃或笛卡尔裂变.
\end{itemize}

    \paragraph{数据分析中的作用与用法}\label{ux6570ux636eux5206ux6790ux4e2dux7684ux4f5cux7528ux4e0eux7528ux6cd5}

    Q2～Q4 基本覆盖了报表分析最常见的三种输出:

\begin{itemize}
\tightlist
\item
  一个按小时/按区域统计的汇总表;
\item
  一张增加了组内特征的新表;
\item
  两张业务表拼成的一张宽表.
\end{itemize}

    \paragraph{Q2:一天中什么时候最忙?}\label{q2ux4e00ux5929ux4e2dux4ec0ux4e48ux65f6ux5019ux6700ux5fd9}

    这个问题的解决思路是: \textgreater{} 先把上车时间拆成小时,
再按小时把订单分组, 对每组统计订单量、总流水和平均客单价;
如果想看更细的模式, 再按``工作日 + 小时''做联合分组.可以按下面顺序思考:

\begin{enumerate}
\def\labelenumi{\arabic{enumi}.}
\tightlist
\item
  用 \texttt{.dt.hour}、\texttt{.dt.day\_name()} 从时间列里提取业务维度;
\item
  用 \texttt{assign} 把这些派生列安全地加回 DataFrame;
\item
  用 \texttt{groupby(...).agg(...)} 输出小时级汇总报表;
\item
  如果要看二维热度分布, 就对多个键一起 \texttt{groupby};
\item
  这一题的核心是 \texttt{groupby\ +\ agg}, \texttt{transform} 和
  \texttt{merge} 在这里不是必须步骤.
\end{enumerate}

    \paragraph{Q3:什么时间段的流水和单位运营时间收入最高?}\label{q3ux4ec0ux4e48ux65f6ux95f4ux6bb5ux7684ux6d41ux6c34ux548cux5355ux4f4dux8fd0ux8425ux65f6ux95f4ux6536ux5165ux6700ux9ad8}

    这个问题相比 Q2 多了一个关键点:

\begin{quote}
\textbf{不是只统计``多少单'', 而是构造新的业务指标后再分组比较.}
\end{quote}

落到 API 上时, 关键不是直接拿原始列做聚合,
而是先把能回答业务问题的指标算出来.

可以按下面顺序思考:

\begin{enumerate}
\def\labelenumi{\arabic{enumi}.}
\tightlist
\item
  用 \texttt{assign} 一次性构造
  \texttt{pickup\_hour}、\texttt{revenue\_per\_minute}、\texttt{fare\_per\_mile}
  等派生指标;
\item
  用 \texttt{groupby("pickup\_hour").agg(...)} 把订单压缩成小时级报表;
\item
  用 \texttt{sort\_values}
  比较哪个小时总流水最高、哪个小时单位运营时间收入最高;
\item
  用 \texttt{transform("mean")} 把组内均值贴回原表,
  再判断单笔订单相对所在小时平均水平是高还是低.
\end{enumerate}

    \paragraph{Q4:哪些上下车区域最热门?}\label{q4ux54eaux4e9bux4e0aux4e0bux8f66ux533aux57dfux6700ux70edux95e8}

    PULocationID:上车区域 ID(Pick-up Location ID, 关联
taxi\_zone\_lookup.csv 中的地理分区与行政区). DOLocationID:下车区域
ID(Drop-off Location ID).

课堂里可以直接按 LocationID 统计热度, 再读取官方
\texttt{taxi\_zone\_lookup.csv}, 把 LocationID 连接成真实的
Borough、Zone 和 service\_zone 信息.

这个问题的业务翻译是:

\begin{quote}
先分别统计上车区域和下车区域各自出现了多少次, 再把两张统计表按
\texttt{LocationID} 拼起来, 得到综合热度表.
\end{quote}

使用这一节 API 解决时, 可以按下面顺序思考:

\begin{enumerate}
\def\labelenumi{\arabic{enumi}.}
\tightlist
\item
  用 \texttt{value\_counts()} 先拿到每个 \texttt{LocationID} 的出现频次;
\item
  用 \texttt{rename\_axis} 和 \texttt{reset\_index}
  把频次统计结果整理成标准 DataFrame;
\item
  用 \texttt{merge} 把频次表和区域维表拼起来, 补充区域标签;
\item
  再用一次 \texttt{merge} 把上车热度和下车热度合成一张综合热度表;
\item
  最后把综合热度表和真实的 taxi zone lookup 维表做 \texttt{merge}, 并用
  \texttt{validate} 显式校验连接关系, 避免拼表时 silently 出错.
\end{enumerate}

    \subsubsection{Pandas: Resample、Rolling 与窗口思维对应着问题
Q5}\label{pandas-resamplerolling-ux4e0eux7a97ux53e3ux601dux7ef4ux5bf9ux5e94ux7740ux95eeux9898-q5}

    哪些订单和时间窗口可能异常?

Q5 的关键词不是``异常值''本身, 而是:

\begin{quote}
\textbf{异常不能只看单条记录, 还要看它在时间窗口中的相对位置.}
这就是窗口思维.
\end{quote}

一笔 80 美元的订单一定异常吗? 不一定. 但如果某个 15
分钟窗口里的订单金额整体突然抬升,
或某一单远高于其所在小时/窗口的分位数阈值, 那就值得关注.

    \paragraph{Resample、Rolling
与窗口思维}\label{resamplerolling-ux4e0eux7a97ux53e3ux601dux7ef4}

    \subparagraph{本质}\label{ux672cux8d28}

    \texttt{resample} 的本质是:

\begin{quote}
\textbf{把时间轴重新切成固定频率的桶, 再在桶内聚合.}
\end{quote}

\texttt{rolling} 的本质是:

\begin{quote}
\textbf{让一个长度固定的滑动窗口沿着时间轴或序列逐步移动,
每移动一步都做一次局部统计.}
\end{quote}

\texttt{transform} 在这里继续扮演``把组级统计量贴回原表''的角色.

    \subparagraph{底层原理}\label{ux5e95ux5c42ux539fux7406}

    这类 API 成立的前提是:

\begin{itemize}
\tightlist
\item
  时间列必须是 \texttt{datetime64{[}ns{]}} 或 \texttt{DatetimeIndex}
\item
  数据最好按时间排序
\item
  Pandas 通过时间桶切分、索引对齐和窗口边界管理来完成批量计算
\end{itemize}

其中:

\begin{itemize}
\tightlist
\item
  \texttt{resample("1h")} 是按小时分桶;
\item
  \texttt{rolling(24)} 是按固定观察点个数滑窗;
\item
  \texttt{rolling("2H")} 则是按时间跨度滑窗.
\end{itemize}

    \subparagraph{数据分析中的作用与用法}\label{ux6570ux636eux5206ux6790ux4e2dux7684ux4f5cux7528ux4e0eux7528ux6cd5}

    在生产环境中, Q5 常见于:

\begin{itemize}
\tightlist
\item
  监控指标异常波动;
\item
  交易金额突增;
\item
  某时间段订单量断崖变化;
\item
  某笔订单相对同时段分布明显偏高.
\end{itemize}

如果把这个问题翻译成 API 的使用路径, 可以理解为:

\begin{quote}
先在单笔订单层面判断异常, 再把订单放回所属时间窗口,
观察它相对同小时和滑动窗口是否异常.
\end{quote}

使用这一节 API 解决时, 可以按下面顺序思考:

\begin{enumerate}
\def\labelenumi{\arabic{enumi}.}
\tightlist
\item
  用 \texttt{assign} 构造 \texttt{revenue\_per\_minute}
  这类更适合做异常判断的指标;
\item
  用 \texttt{quantile} 和 \texttt{np.select}
  先做单笔订单的分位数异常分类;
\item
  用 \texttt{transform} 计算``所在小时''的中位数和 p95,
  并把它们回填到每一行;
\item
  用 \texttt{where} 只保留异常值, 方便快速抽查;
\item
  用 \texttt{resample} 做小时级窗口汇总, 再用 \texttt{rolling}
  构造动态阈值.
\end{enumerate}

    \subsubsection{Q6:
为什么同一个分析可以相差一个数量级甚至更多?}\label{q6-ux4e3aux4ec0ux4e48ux540cux4e00ux4e2aux5206ux6790ux53efux4ee5ux76f8ux5deeux4e00ux4e2aux6570ux91cfux7ea7ux751aux81f3ux66f4ux591a}

    这个问题的核心不是``谁会写更多 API'', 而是:
\textbf{不同写法把计算放在了不同的执行层.}

    \paragraph{Pandas
密集计算的几种执行方式}\label{pandas-ux5bc6ux96c6ux8ba1ux7b97ux7684ux51e0ux79cdux6267ux884cux65b9ux5f0f}

    \begin{enumerate}
\def\labelenumi{\arabic{enumi}.}
\item
  \texttt{Python\ loop} 最直接, 但循环控制、分支判断、类型解析都发生在
  Python 解释器层, 开销最高.
\item
  \texttt{apply(axis=1)} 看起来像 Pandas 风格, 但本质仍是``逐行调用
  Python 函数''; 每一行都可能被包装成临时 \texttt{Series},
  因此很多场景甚至比手写循环还慢.
\item
  \texttt{NumPy\ /\ Pandas\ vectorization} 把整列计算下沉到底层数组和
  ufunc, 是数值计算的默认首选.只要表达式足够简单、分支不多,
  通常都能跑得很快.
\item
  \texttt{Numba} 当计算已经超出``自然向量化''的舒适区,
  例如有多层分支、反复更新、临时量较多时, 可以把数值循环 JIT 成机器码,
  常常会比纯向量化更快.
\end{enumerate}

    \paragraph{Pandas
的内存与存储优化}\label{pandas-ux7684ux5185ux5b58ux4e0eux5b58ux50a8ux4f18ux5316}

    \begin{enumerate}
\def\labelenumi{\arabic{enumi}.}
\tightlist
\item
  合适的 \texttt{dtype} 会直接影响内存占用、缓存命中率和整体吞吐.
\item
  性能实验前应尽早裁剪列, 只读取真正要用的字段.
\item
  \texttt{Parquet} 是带 schema 的列式二进制格式, 通常比 \texttt{CSV}
  更适合分析场景.
\item
  做性能比较时, 应该尽量把 ``I/O / 清洗 / 核心计算'' 三件事拆开,
  不要把它们混在一个 benchmark 里.
\end{enumerate}

    \subsubsection{这一部分 Pandas
总结}\label{ux8fd9ux4e00ux90e8ux5206-pandas-ux603bux7ed3}

    如果把 Pandas 这一大段内容压缩成一句话, 那就是:

\begin{quote}
\textbf{Pandas 不是``带表头的 Excel 替代品'',
而是一套把标签、缺失值、分组、连接、时间窗口和性能工程统一起来的表格分析计算模型.}
\end{quote}

从 Q1 到 Q6, 学员真正应该建立起来的不是零散 API 记忆,
而是下面这条分析路径:

\begin{Shaded}
\begin{Highlighting}[]
\NormalTok{Q1 数据是否可信? }
\NormalTok{    {-}\textgreater{} dtype / missing / mask / loc / query}

\NormalTok{Q2 指标怎么聚合? }
\NormalTok{    {-}\textgreater{} datetime / groupby / agg}

\NormalTok{Q3 聚合结果怎么回填成特征? }
\NormalTok{    {-}\textgreater{} assign / transform / 派生指标}

\NormalTok{Q4 多张表怎么安全拼接? }
\NormalTok{    {-}\textgreater{} value\_counts / merge / validate}

\NormalTok{Q5 如何引入时间窗口思维? }
\NormalTok{    {-}\textgreater{} quantile / where / resample / rolling}

\NormalTok{Q6 为什么性能差异会这么大? }
\NormalTok{    {-}\textgreater{} loop / apply / NumPy / dtype / query/eval / Numba / Parquet}
\end{Highlighting}
\end{Shaded}

    \subsection{Python
数据分析的扩展}\label{python-ux6570ux636eux5206ux6790ux7684ux6269ux5c55}

    前面的 Q1～Q6 主要解决的是:

\begin{itemize}
\tightlist
\item
  如何理解数据;\\
\item
  如何清洗数据;\\
\item
  如何做分组、窗口和性能分析;\\
\item
  如何在 Pandas / NumPy 里把分析问题写清楚.
\end{itemize}

但在真实项目里, 课程通常还要再往前走一步:

\begin{enumerate}
\def\labelenumi{\arabic{enumi}.}
\tightlist
\item
  把一组稳定重复的分析动作沉淀成\textbf{领域 API},
  避免团队成员反复复制粘贴脚本;
\item
  把真正的热点数值循环下沉到\textbf{原生编译层}, 避免所有性能问题都卡在
  Python 解释器上.
\end{enumerate}

下面用两个非常典型的扩展方向来收尾:

\begin{enumerate}
\def\labelenumi{\arabic{enumi}.}
\tightlist
\item
  \texttt{DataFrame\ Accessor}:
  解决``怎么把分析脚本升级成可复用领域工具'';
\item
  \texttt{Numba\ JIT}: 解决``怎么把复杂数值循环编译成机器码''.
\end{enumerate}

    \subsubsection{自定义
Accessor(描述符模式与命名空间扩展)}\label{ux81eaux5b9aux4e49-accessorux63cfux8ff0ux7b26ux6a21ux5f0fux4e0eux547dux540dux7a7aux95f4ux6269ux5c55}

    \paragraph{原理}\label{ux539fux7406}

    Accessor 可以理解为: \textbf{在不修改 Pandas 源码的前提下, 给
\texttt{DataFrame} 注册一个自定义业务命名空间}.

例如, 当我们写:

\begin{Shaded}
\begin{Highlighting}[]
\NormalTok{df.taxi.clean()}
\end{Highlighting}
\end{Shaded}

这里的 \texttt{taxi} 不是 DataFrame 的原生字段, 也不是普通列名,
而是我们通过:

\begin{Shaded}
\begin{Highlighting}[]
\AttributeTok{@pd.api.extensions.register\_dataframe\_accessor}\NormalTok{(}\StringTok{"taxi"}\NormalTok{)}
\end{Highlighting}
\end{Shaded}

注册出来的一个领域入口.

从使用体验上看, 它很像一种``描述符式的命名空间扩展'':

\begin{enumerate}
\def\labelenumi{\arabic{enumi}.}
\tightlist
\item
  Pandas 在 \texttt{DataFrame} 类上挂上一个叫 \texttt{taxi} 的访问入口;
\item
  当你访问 \texttt{df.taxi} 时, Pandas 会创建一个和当前 \texttt{df}
  绑定的 accessor 对象;
\item
  accessor 内部通常把原始 \texttt{DataFrame} 保存在 \texttt{self.\_obj}
  中;
\item
  然后你就可以在这个对象上暴露
  \texttt{clean()}、\texttt{hourly\_kpis()}、\texttt{location\_popularity()}
  这类业务方法.
\end{enumerate}

它的价值不在于``语法更酷'', 而在于:

\begin{enumerate}
\def\labelenumi{\arabic{enumi}.}
\tightlist
\item
  把零散脚本沉淀成领域 API;
\item
  让团队成员直接复用统一规则;
\item
  把字段校验、特征工程、清洗逻辑集中到一个地方维护;
\item
  让后续代码从``操作表''升级为``调用业务能力''.
\end{enumerate}

    \paragraph{使用场景}\label{ux4f7fux7528ux573aux666f}

    Accessor 特别适合下面这类情况:

\begin{enumerate}
\def\labelenumi{\arabic{enumi}.}
\tightlist
\item
  同一份业务表反复做相同的字段校验、特征构造和清洗规则;
\item
  团队里已经形成稳定的分析词汇, 比如``可信订单''\,``小时
  KPI''\,``区域热度'';
\item
  希望下游代码更接近业务表达, 而不是每次都重新手搓几十行 Pandas;
\item
  希望把分析脚本逐步升级成公司内部可复用的工具层.
\end{enumerate}

    \subsubsection{Numba JIT 加速与绕过 GIL(Native
Compilation)}\label{numba-jit-ux52a0ux901fux4e0eux7ed5ux8fc7-gilnative-compilation}

    \paragraph{原理}\label{ux539fux7406}

    Numba 的核心不是``把 Pandas 变快'', 而是:
\textbf{把已经整理好的数值循环编译成原生机器码}.

也就是说, Numba 最擅长的不是下面这些:

\begin{enumerate}
\def\labelenumi{\arabic{enumi}.}
\tightlist
\item
  读 CSV / Parquet;
\item
  做 \texttt{merge};
\item
  做 \texttt{groupby};
\item
  处理大量字符串、对象列和标签对齐.
\end{enumerate}

这些事情仍然更适合交给 Pandas.

Numba 真正擅长的是:

\begin{enumerate}
\def\labelenumi{\arabic{enumi}.}
\tightlist
\item
  输入已经是 NumPy 数组;
\item
  计算逻辑里有大量数值循环;
\item
  中间有多层 \texttt{if/else} 分支;
\item
  重复更新很多次, 用纯向量化会写得非常绕, 或者会制造很多中间临时数组.
\end{enumerate}

典型写法是:

\begin{Shaded}
\begin{Highlighting}[]
\AttributeTok{@njit}
\KeywordTok{def}\NormalTok{ kernel(...):}
\NormalTok{    ...}
\end{Highlighting}
\end{Shaded}

它的原理可以概括成三句话:

\begin{enumerate}
\def\labelenumi{\arabic{enumi}.}
\tightlist
\item
  \texttt{@njit} 会让 Numba 尝试进入 \texttt{nopython} 模式,
  把循环编译成原生代码;
\item
  编译成功后, 循环中的数值运算不再依赖 Python 解释器逐条调度;
\item
  如果再加上 \texttt{nogil=True}, 那么进入这段编译后代码时,
  线程可以不一直持有 GIL; 如果加上 \texttt{parallel=True} 和
  \texttt{prange}, 还可以在满足独立迭代的前提下做多线程并行.
\end{enumerate}

要特别强调一点:

\begin{quote}
\textbf{Numba 不是用来替代 Pandas 的, 它更像是``把 Pandas 清洗好的数组,
交给原生数值内核去算''.}
\end{quote}

    \paragraph{使用场景}\label{ux4f7fux7528ux573aux666f}

    Numba 特别适合下面这类情况:

\begin{enumerate}
\def\labelenumi{\arabic{enumi}.}
\tightlist
\item
  行数很多, 且每行都要做复杂数值计算;
\item
  逻辑中有多层分支、重复更新、局部状态累积;
\item
  用纯向量化虽然能写, 但会出现很多 \texttt{np.where}、很多中间数组,
  代码可读性和内存压力都很差;
\item
  每一行之间相互独立, 适合 \texttt{parallel=True} + \texttt{prange}
  并行.
\end{enumerate}

    \begin{tcolorbox}[breakable, size=fbox, boxrule=1pt, pad at break*=1mm,colback=cellbackground, colframe=cellborder]
\prompt{In}{incolor}{ }{\boxspacing}
\begin{Verbatim}[commandchars=\\\{\}]

\end{Verbatim}
\end{tcolorbox}


    % Add a bibliography block to the postdoc
    
    
    
\end{document}
